%%%
%%% Radical "⼼"
%%%
\section*{Radical 61: ``⼼'' (忄、⺗)}\addcontentsline{toc}{section}{Radical 61: ⼼,忄、⺗}\addcontentsline{loh}{figure}{\#\#\#\# 61: ⼼}

%%%%%%%%%% 心 %%%%%%%%%%
\subsection*{心}\addcontentsline{loh}{figure}{心}

\begin{Entry}{心}{4}{⼼}[Kangxi 61]
  \begin{Phonetics}{心}{xin1}[][HSK 3]
    \definition*{s.}{Xin, uma das mansões lunares; uma das vinte e oito constelações}
    \definition[颗,个]{s.}{o coraçã; órgão que impulsiona a circulação sanguínea no corpo humano e nos vertebrados | coração; mente; sentimento; intenção; refere"-se aos órgãos do pensamento e ao pensamento, sentimentos, etc. | centro; núcleo; parte central}
  \end{Phonetics}
\end{Entry}

\begin{Entry}{心中}{4,4}{⼼,⼁}
  \begin{Phonetics}{心中}{xin1zhong1}[][HSK 2]
    \definition{s.}{no coração; na mente}
  \synonymref{内心}{nei4xin1}
  \synonymref{心里}{xin1li3}
  \end{Phonetics}
\end{Entry}

\begin{Entry}{心目}{4,5}{⼼,⽬}
  \begin{Phonetics}{心目}{xin1mu4}[][HSK 7-9]
    \definition{s.}{humor; estado de espírito; refere"-se a sentimentos na mente ou nos sentidos visuais | mente; visão mental; refere"-se a pensamentos e opiniões | memória}
  \synonymref{内心}{nei4xin1}
  \end{Phonetics}
\end{Entry}

\begin{Entry}{心安}{4,6}{⼼,⼧}
  \begin{Phonetics}{心安}{xin1'an1}
    \definition{adj.}{à vontade; tranquilizado}
  \seealsoref{安心}{an1xin1}
  \end{Phonetics}
\end{Entry}

\begin{Entry}{心安理得}{4,6,11,11}{⼼,⼧,⽟,⼻}
  \begin{Phonetics}{心安理得}{xin1'an1-li3de2}[][HSK 7-9]
    \definition{expr.}{sentir"-se à vontade e justificado; ter a consciência tranquila; com a mente em paz e a consciência limpa; ter a consciência limpa; não ter receio algum de algo}
  \synonymref{理直气壮}{li3zhi2-qi4zhuang4}
  \antonymref{提心吊胆}{ti2xin1-diao4dan3}
  \end{Phonetics}
\end{Entry}

\begin{Entry}{心机}{4,6}{⼼,⽊}
  \begin{Phonetics}{心机}{xin1ji1}
    \definition{s.}{ideias; estratégias; pensamentos; planos; tramas}
  \synonymref{心思}{xin1si5}
  \antonymref{心绪}{xin1xu4}
  \end{Phonetics}
\end{Entry}

\begin{Entry}{心血}{4,6}{⼼,⾎}
  \begin{Phonetics}{心血}{xin1xue4}[][HSK 7-9]
    \definition{s.}{esforço meticuloso; esforço empregado para alcançar o sucesso}
  \synonymref{血汗}{xue4han4}
  \end{Phonetics}
\end{Entry}

\begin{Entry}{心声}{4,7}{⼼,⼠}
  \begin{Phonetics}{心声}{xin1sheng1}[][HSK 7-9]
    \definition{s.}{desejos sinceros; aspirações; pensamentos}
  \synonymref{渴望}{ke3wang4}
  \synonymref{心事}{xin1shi4}
  \end{Phonetics}
\end{Entry}

\begin{Entry}{心灵}{4,7}{⼼,⽕}
  \begin{Phonetics}{心灵}{xin1ling2}[][HSK 6]
    \definition[个,颗]{s.}{alma; coração; espírito; refere"-se ao coração, espírito, pensamentos, etc.}
  \end{Phonetics}
\end{Entry}

\begin{Entry}{心灵手巧}{4,7,4,5}{⼼,⽕,⼿,⼯}
  \begin{Phonetics}{心灵手巧}{xin1ling2-shou3qiao3}[][HSK 7-9]
    \definition{expr.}{inteligente e habilidoso; capaz; hábil}
  \end{Phonetics}
\end{Entry}

\begin{Entry}{心肠}{4,7}{⼼,⾁}
  \begin{Phonetics}{心肠}{xin1chang2}[][HSK 7-9]
    \definition{s.}{coração; intenção; estado emocional | preocupações; coisas que pesam na mente; pensamentos; assuntos do coração;  refere"-se a sentimentos pessoais, inquietações ou coisas que perturbam a mente | humor}
  \end{Phonetics}
\end{Entry}

\begin{Entry}{心里}{4,7}{⼼,⾥}
  \begin{Phonetics}{心里}{xin1li3}[][HSK 2]
    \definition[个]{s.}{no coração; no coração de alguém | no coração; na mente; na cabeça e no peito}
  \synonymref{内心}{nei4xin1}
  \synonymref{心胸}{xin1xiong1}
  \synonymref{心中}{xin1zhong1}
  \synonymref{胸膛}{xiong1tang2}
  \antonymref{表面}{biao3mian4}
  \antonymref{外表}{wai4biao3}
  \end{Phonetics}
\end{Entry}

\begin{Entry}{心里话}{4,7,8}{⼼,⾥,⾔}
  \begin{Phonetics}{心里话}{xin1li3hua4}[][HSK 7-9]
    \definition[句]{s.}{os pensamentos e sentimentos mais íntimos de alguém}
  \end{Phonetics}
\end{Entry}

\begin{Entry}{心事}{4,8}{⼼,⼅}
  \begin{Phonetics}{心事}{xin1shi4}[][HSK 7-9]
    \definition[桩,件]{s.}{mágoa; preocupação; algo que pesa na mente; um fardo para a mente; as coisas em que estou pensando (geralmente me referindo a coisas que me deixam desconfortável ou constrangido)}
  \synonymref{心思}{xin1si5}
  \end{Phonetics}
\end{Entry}

\begin{Entry}{心态}{4,8}{⼼,⼼}
  \begin{Phonetics}{心态}{xin1tai4}[][HSK 5]
    \definition[种,个]{s.}{mentalidade; psicologia; estado mental}
  \end{Phonetics}
\end{Entry}

\begin{Entry}{心思}{4,9}{⼼,⼼}
  \begin{Phonetics}{心思}{xin1si5}[][HSK 7-9]
    \definition[番]{s.}{pensamento; ideia | pensamento; refere"-se a habilidades como raciocínio e memória | humor; estado de espírito; a sensação de querer fazer algo}
  \seealsoref{心事}{xin1shi4}
  \synonymref{脑筋}{nao3jin1}
  \synonymref{情绪}{qing2xu4}
  \synonymref{思想}{si1xiang3}
  \synonymref{头脑}{tou2nao3}
  \synonymref{心情}{xin1qing2}
  \synonymref{心绪}{xin1xu4}
  \end{Phonetics}
\end{Entry}

\begin{Entry}{心急如火}{4,9,6,4}{⼼,⼼,⼥,⽕}
  \begin{Phonetics}{心急如火}{xin1ji2-ru2huo3}
    \definition{expr.}{ansioso e impaciente; ardendo de impaciência}
  \end{Phonetics}
\end{Entry}

\begin{Entry}{心急如焚}{4,9,6,12}{⼼,⼼,⼥,⽕}
  \begin{Phonetics}{心急如焚}{xin1ji2-ru2fen2}[][HSK 7-9]
    \definition{s.}{ansioso; ardendo de impaciência}
  \seealsoref{心急如火}{xin1ji2-ru2huo3}
  \end{Phonetics}
\end{Entry}

\begin{Entry}{心爱}{4,10}{⼼,⽖}
  \begin{Phonetics}{心爱}{xin1'ai4}[][HSK 7-9]
    \definition{v.}{amar; valorizar; ter afeição sincera; ter um relacionamento próximo; ter afeto profundo}
  \synonymref{可爱}{ke3'ai4}
  \synonymref{亲爱}{qin1'ai4}
  \synonymref{热爱}{re4'ai4}
  \synonymref{喜爱}{xi3'ai4}
  \synonymref{喜欢}{xi3huan5}
  \end{Phonetics}
\end{Entry}

\begin{Entry}{心疼}{4,10}{⼼,⽧}
  \begin{Phonetics}{心疼}{xin1teng2}[][HSK 5]
    \definition{v.}{amar profundamente; sentir pena porque coisas valiosas foram destruídas ou perdidas; não querer se separar delas | sentir pena; ficar angustiado; preocupar"-se e sofrer pelo sofrimento dos outros; estar disposto a cuidar mais por causa da preocupação}
  \end{Phonetics}
\end{Entry}

\begin{Entry}{心病}{4,10}{⼼,⽧}
  \begin{Phonetics}{心病}{xin1bing4}[][HSK 7-9]
    \definition{s.}{preocupação; ansiedade; humor deprimido | ponto sensível; problema secreto; histórias ocultas ou dor oculta | doença cardíaca}
  \end{Phonetics}
\end{Entry}

\begin{Entry}{心胸}{4,10}{⼼,⾁}
  \begin{Phonetics}{心胸}{xin1xiong1}[][HSK 7-9]
    \definition{s.}{mentalidade; tolerância; mente aberta; amplitude de espírito | aspiração; ambição}
  \seealsoref{胸怀}{xiong1huai2}
  \end{Phonetics}
\end{Entry}

\begin{Entry}{心脏}{4,10}{⼼,⾁}
  \begin{Phonetics}{心脏}{xin1zang4}[][HSK 6]
    \definition[颗,个]{s.}{coração; um órgão importante no corpo de humanos ou animais superiores que faz o sangue circular | coração; o centro ou a parte mais importante de uma metáfora}
  \end{Phonetics}
\end{Entry}

\begin{Entry}{心脏病}{4,10,10}{⼼,⾁,⽧}
  \begin{Phonetics}{心脏病}{xin1zang4bing4}[][HSK 6]
    \definition{s.}{doença cardíaca; cardiopatia um termo geral para anormalidades ou doenças na estrutura e função do coração humano}
  \end{Phonetics}
\end{Entry}

\begin{Entry}{心得}{4,11}{⼼,⼻}
  \begin{Phonetics}{心得}{xin1de2}[][HSK 7-9]
    \definition[篇,点,项,个]{s.}{percepção; experiência; conhecimento adquirido; os princípios que se aprendem e se vivenciam no trabalho e nos estudos}
  \end{Phonetics}
\end{Entry}

\begin{Entry}{心情}{4,11}{⼼,⼼}
  \begin{Phonetics}{心情}{xin1qing2}[][HSK 2]
    \definition{s.}{humor; tom de sentimento; estado de espírito; estado emocional interior}
  \seealsoref{情绪}{qing2xu4}
  \seealsoref{心绪}{xin1xu4}
  \synonymref{感情}{gan3qing2}
  \synonymref{情感}{qing2gan3}
  \synonymref{情绪}{qing2xu4}
  \synonymref{心理}{xin1li3}
  \synonymref{心绪}{xin1xu4}
  \synonymref{心思}{xin1si5}
  \end{Phonetics}
\end{Entry}

\begin{Entry}{心理}{4,11}{⼼,⽟}
  \begin{Phonetics}{心理}{xin1li3}[][HSK 4]
    \definition[个]{s.}{mentalidade; refere"-se à reflexão da mente humana sobre coisas objetivas, incluindo sensação, percepção, memória, pensamento e emoções | psicologia}
  \end{Phonetics}
\end{Entry}

\begin{Entry}{心眼儿}{4,11,2}{⼼,⽬,⼉}
  \begin{Phonetics}{心眼儿}{xin1yan3r5}[][HSK 7-9]
    \definition{s.}{coração; mente; intenção; a mente de uma pessoa | inteligência; astúcia; engenhosidade | dúvidas infundadas; receios desnecessários; preocupações e considerações desnecessárias sobre as pessoas | tolerância}
  \end{Phonetics}
\end{Entry}

\begin{Entry}{心绪}{4,11}{⼼,⽷}
  \begin{Phonetics}{心绪}{xin1xu4}
    \definition{s.}{estado de espírito | sentimentos; mente; humor (geralmente descrito como estável ou desordenado)}
  \seealsoref{心情}{xin1qing2}
  \synonymref{情绪}{qing2xu4}
  \synonymref{心理}{xin1li3}
  \synonymref{心情}{xin1qing2}
  \synonymref{心思}{xin1si5}
  \end{Phonetics}
\end{Entry}

\begin{Entry}{心慌}{4,12}{⼼,⼼}
  \begin{Phonetics}{心慌}{xin1/huang1}[][HSK 7-9]
    \definition{adj.}{perturbado; nervoso; alarmado; ansioso}
    \definition{v.+compl.}{ter palpitações; ter aumento ou irregularidade na frequência cardíaca}
  \synonymref{慌乱}{huang1luan4}
  \synonymref{慌张}{huang1zhang1}
  \end{Phonetics}
\end{Entry}

\begin{Entry}{心想事成}{4,13,8,6}{⼼,⼼,⼅,⼽}
  \begin{Phonetics}{心想事成}{xin1xiang3-shi4cheng2}[][HSK 7-9]
    \definition{expr.}{``Que todos os seus desejos se realizem.''; ter os próprios desejos realizados; independentemente do que você imaginar, você terá sucesso}
  \synonymref{如愿以偿}{ru2yuan4yi3chang2}
  \synonymref{一帆风顺}{yi4fan1-feng1shun4}
  \end{Phonetics}
\end{Entry}

\begin{Entry}{心意}{4,13}{⼼,⼼}
  \begin{Phonetics}{心意}{xin1yi4}[][HSK 7-9]
    \definition{s.}{consideração; propósito; sentimentos amáveis; os sentimentos positivos em relação aos outros geralmente são expressos por meio de ajuda, agradecimento e votos de felicidades | mente; o que eu quero expressar; meus pensamentos mais íntimos}
  \seealsoref{心愿}{xin1yuan4}
  \end{Phonetics}
\end{Entry}

\begin{Entry}{心愿}{4,14}{⼼,⽕}
  \begin{Phonetics}{心愿}{xin1yuan4}[][HSK 6]
    \definition[桩]{s.}{desejo acalentado; aspiração; desejo; sonho | o desejo do coração}
  \end{Phonetics}
\end{Entry}

\begin{Entry}{心酸}{4,14}{⼼,⾣}
  \begin{Phonetics}{心酸}{xin1suan1}[][HSK 7-9]
    \definition{adj.}{de partir o coração; triste; dor de coração causada pelo luto}
  \synonymref{悲哀}{bei1'ai1}
  \synonymref{悲伤}{bei1shang1}
  \synonymref{辛酸}{xin1suan1}
  \antonymref{甜蜜}{tian2mi4}
  \end{Phonetics}
\end{Entry}

%%%%%%%%%% 必 %%%%%%%%%%
\subsection*{必}\addcontentsline{loh}{figure}{必}

\begin{Entry}{必}{5}{⼼}
  \begin{Phonetics}{必}{bi4}[][HSK 5]
    \definition{adv.}{certamente; necessariamente; indica que algo é certo ou que alguém acredita que esteja correto | deve; tem que}
  \seealsoref{必须}{bi4xu1}
  \end{Phonetics}
\end{Entry}

\begin{Entry}{必不可少}{5,4,5,4}{⼼,⼀,⼝,⼩}
  \begin{Phonetics}{必不可少}{bi4bu4ke3shao3}[][HSK 7-9]
    \definition{suf.}{indispensável; essencial; absolutamente necessário; insubstituível; inevitável}
  \end{Phonetics}
\end{Entry}

\begin{Entry}{必定}{5,8}{⼼,⼧}
  \begin{Phonetics}{必定}{bi4ding4}[][HSK 7-9]
    \definition{adv.}{certamente; estar vinculado a; ter certeza de; para expressar a certeza de um julgamento ou inferência | definitivamente; para expressar determinação de vontade; ter certeza de fazê"-lo}
  \seealsoref{一定}{yi2ding4}
  \synonymref{必将}{bi4jiang1}
  \synonymref{必然}{bi4ran2}
  \synonymref{必需}{bi4xu1}
  \synonymref{肯定}{ken3ding4}
  \synonymref{一定}{yi2ding4}
  \synonymref{注定}{zhu4ding4}
  \antonymref{不必}{bu2bi4}
  \antonymref{恐怕}{kong3pa4}
  \antonymref{也许}{ye3xu3}
  \end{Phonetics}
\end{Entry}

\begin{Entry}{必修}{5,9}{⼼,⼈}
  \begin{Phonetics}{必修}{bi4xiu1}[][HSK 6]
    \definition{adj.}{(de um curso acadêmico) obrigatório; compulsório; mandatório; obrigatório estudar de acordo com os regulamentos}
  \antonymref{选修}{xuan3xiu1}
  \end{Phonetics}
\end{Entry}

\begin{Entry}{必将}{5,9}{⼼,⼨}
  \begin{Phonetics}{必将}{bi4jiang1}[][HSK 6]
    \definition{adv.}{certamente; certamente irá; usado para expressar inevitabilidade (ou necessidade)}
  \seealsoref{将}{jiang1}
  \synonymref{必定}{bi4ding4}
  \synonymref{必然}{bi4ran2}
  \synonymref{即将}{ji2jiang1}
  \antonymref{未必}{wei4bi4}
  \end{Phonetics}
\end{Entry}

\begin{Entry}{必要}{5,9}{⼼,⾑}
  \begin{Phonetics}{必要}{bi4yao4}[][HSK 3]
    \definition{adj.}{necessário; essencial; indispensável}
    \definition[个,些]{s.}{necessidade; características indispensáveis}
  \seealsoref{必须}{bi4xu1}
  \synonymref{必需}{bi4xu1}
  \synonymref{必须}{bi4xu1}
  \synonymref{需求}{xu1qiu2}
  \synonymref{需要}{xu1yao4}
  \antonymref{不必}{bu2bi4}
  \antonymref{不要}{bu2yao4}
  \antonymref{多余}{duo1yu2}
  \antonymref{无须}{wu2xu1}
  \end{Phonetics}
\end{Entry}

\begin{Entry}{必须}{5,9}{⼼,⾴}
  \begin{Phonetics}{必须}{bi4xu1}[][HSK 2]
    \definition{adv.}{necessariamente; obrigatoriamente; indica a necessidade lógica e emocional | deve; tem que; é obrigado a}
  \seealsoref{必}{bi4}
  \seealsoref{必需}{bi4xu1}
  \synonymref{必}{bi4}
  \synonymref{必需}{bi4xu1}
  \synonymref{必要}{bi4yao4}
  \synonymref{不得不}{bu4de2bu4}
  \synonymref{得}{dei3}
  \synonymref{非得}{fei1dei3}
  \synonymref{务必}{wu4bi4}
  \synonymref{一定}{yi2ding4}
  \antonymref{不必}{bu2bi4}
  \antonymref{不用}{bu2yong4}
  \antonymref{无须}{wu2xu1}
  \end{Phonetics}
\end{Entry}

\begin{Entry}{必然}{5,12}{⼼,⽕}
  \begin{Phonetics}{必然}{bi4ran2}[][HSK 3]
    \definition{adj.}{certo; inevitável; necessário; definido e inalterável; imutável}
    \definition{adv.}{inevitavelmente}
    \definition{s.}{necessidade; em filosofia, refere"-se às leis objetivas do desenvolvimento que não são influenciadas pela vontade humana}
  \seealsoref{一定}{yi2ding4}
  \seealsoref{自然}{zi4ran5}
  \synonymref{必定}{bi4ding4}
  \synonymref{必将}{bi4jiang1}
  \synonymref{当然}{dang1ran2}
  \synonymref{绝对}{jue2dui4}
  \synonymref{肯定}{ken3ding4}
  \synonymref{势必}{shi4bi4}
  \synonymref{一定}{yi2ding4}
  \synonymref{自然}{zi4ran5}
  \antonymref{或者}{huo4zhe3}
  \antonymref{绝不}{jue2bu4}
  \antonymref{恐怕}{kong3pa4}
  \antonymref{偶然}{ou3ran2}
  \antonymref{碰巧}{peng4qiao3}
  \antonymref{未必}{wei4bi4}
  \antonymref{也许}{ye3xu3}
  \end{Phonetics}
\end{Entry}

\begin{Entry}{必需}{5,14}{⼼,⾬}
  \begin{Phonetics}{必需}{bi4xu1}[][HSK 5]
    \definition{adj.}{essencial; indispensável}
    \definition{v.}{ser essencial; ser indispensável}
  \seealsoref{必须}{bi4xu1}
  \synonymref{必定}{bi4ding4}
  \synonymref{必须}{bi4xu1}
  \synonymref{必要}{bi4yao4}
  \synonymref{务必}{wu4bi4}
  \synonymref{一定}{yi2ding4}
  \antonymref{不必}{bu2bi4}
  \end{Phonetics}
\end{Entry}

%%%%%%%%%% 忙 %%%%%%%%%%
\subsection*{忙}\addcontentsline{loh}{figure}{忙}

\begin{Entry}{忙}{6}{⼼}
  \begin{Phonetics}{忙}{mang2}[][HSK 1]
    \definition*{s.}{Sobrenome: Mang}
    \definition{adj.}{ocupado; movimentado; totalmente ocupado; muitas coisas para fazer, sem tempo livre | imperativo; ansioso; urgente}
    \definition{v.}{apressar"-se; agitar"-se; fazer algo com urgência e constantemente | trabalhar; fazer}
  \antonymref{闲}{xian2}
  \end{Phonetics}
\end{Entry}

\begin{Entry}{忙乱}{6,7}{⼼,⼄}
  \begin{Phonetics}{忙乱}{mang2luan4}[][HSK 7-9]
    \definition{adj.}{apressado e desordenado; às pressas e em meio à confusão}
    \definition{v.}{estar com pressa e desorganizado; realizar uma tarefa de forma apressada e desordenada; agir de forma apressada e desordenada}
  \end{Phonetics}
\end{Entry}

\begin{Entry}{忙活}{6,9}{⼼,⽔}
  \begin{Phonetics}{忙活}{mang2huo2}
    \definition{s.}{tarefa urgente}
    \definition{v.}{estar ocupado com o trabalho; estar envolvido com tarefas; estar com pressa para terminar as coisas}
  \end{Phonetics}
  \begin{Phonetics}{忙活}{mang2huo5}[][HSK 7-9]
    \definition{adj.}{Coloquial: ocupado; movimentado}
    \definition{v.}{estar ocupado; movimentar"-se freneticamente}
  \end{Phonetics}
\end{Entry}

\begin{Entry}{忙得}{6,11}{⼼,⼻}
  \begin{Phonetics}{忙得}{mang2 de2}
    \definition{adj.}{muito ocupado}
  \end{Phonetics}
\end{Entry}

\begin{Entry}{忙碌}{6,13}{⼼,⽯}
  \begin{Phonetics}{忙碌}{mang2lu4}[][HSK 7-9]
    \definition{adj.}{ocupado; atarefado; agitado; movimentado; ocupado com várias coisas, sem tempo livre}
  \end{Phonetics}
\end{Entry}

%%%%%%%%%% 忌 %%%%%%%%%%
\subsection*{忌}\addcontentsline{loh}{figure}{忌}

\begin{Entry}{忌}{7}{⼼}
  \begin{Phonetics}{忌}{ji4}[][HSK 7-9]
    \definition{s.}{medo; pavor; escrúpulo}
    \definition{v.}{ter ciúmes de; invejar | evitar; afastar"-se de; esquivar"-se de ; abster"-se de | desistir; desistir}
  \end{Phonetics}
\end{Entry}

\begin{Entry}{忌口}{7,3}{⼼,⼝}
  \begin{Phonetics}{忌口}{ji4/kou3}[][HSK 7-9]
    \definition{v.+compl.}{evitar certos alimentos (como quando se está doente); estar de dieta}
  \end{Phonetics}
\end{Entry}

\begin{Entry}{忌讳}{7,6}{⼼,⾔}
  \begin{Phonetics}{忌讳}{ji4hui4}[][HSK 7-9]
    \definition[个,种]{s.}{tabu; certas palavras ou ações são tabu devido a costumes ou razões pessoais}
    \definition{v.}{evitar; ser tabu sobre; tentar evitar ou não querer que aconteça (algo que pode ter consequências negativas)}
  \end{Phonetics}
\end{Entry}

\begin{Entry}{忌妒}{7,7}{⼼,⼥}
  \begin{Phonetics}{忌妒}{ji4du4}
    \definition{v.}{invejar; ter ciúmes de; estar infeliz ou até mesmo odiar ou ter ciúmes dos outros porque eles são melhores que eles em termos de talento, status, etc.}
  \end{Phonetics}
\end{Entry}

%%%%%%%%%% 忍 %%%%%%%%%%
\subsection*{忍}\addcontentsline{loh}{figure}{忍}

\begin{Entry}{忍}{7}{⼼}
  \begin{Phonetics}{忍}{ren3}[][HSK 5]
    \definition{v.}{suportar; aguentar; tolerar; aturar | ter coragem para; ser insensível o suficiente para; ser capaz de endurecer o coração e fazer coisas que não se devem fazer por uma questão de razão}
  \end{Phonetics}
\end{Entry}

\begin{Entry}{忍不住}{7,4,7}{⼼,⼀,⼈}
  \begin{Phonetics}{忍不住}{ren3bu5zhu4}[][HSK 5]
    \definition{v.}{incapaz de suportar; não conseguir evitar fazer algo; não conseguir se controlar}
  \end{Phonetics}
\end{Entry}

\begin{Entry}{忍心}{7,4}{⼼,⼼}
  \begin{Phonetics}{忍心}{ren3/xin1}[][HSK 7-9]
    \definition{v.+compl.}{ter a coragem de; ser suficientemente insensível para; ser capaz de endurecer o coração (fazer coisas que não se suporta fazer)}
  \end{Phonetics}
\end{Entry}

\begin{Entry}{忍饥挨饿}{7,5,10,10}{⼼,⾷,⼿,⾷}
  \begin{Phonetics}{忍饥挨饿}{ren3ji1-ai2'e4}[][HSK 7-9]
    \definition{expr.}{``Morrendo de fome.''; suportar os tormentos da fome; famintos}
  \end{Phonetics}
\end{Entry}

\begin{Entry}{忍受}{7,8}{⼼,⼜}
  \begin{Phonetics}{忍受}{ren3shou4}[][HSK 5]
    \definition{v.}{suportar; sofrer; aguentar; tolerar; suportar com dificuldade o sofrimento, as dificuldades e as adversidades da vida}
  \end{Phonetics}
\end{Entry}

\begin{Entry}{忍耐}{7,9}{⼼,⽽}
  \begin{Phonetics}{忍耐}{ren3nai4}[][HSK 7-9]
    \definition{v.}{conter"-se; exercer paciência (ou contenção); suprimir um determinado sentimento ou emoção para impedir que ele seja expresso}
  \end{Phonetics}
\end{Entry}

%%%%%%%%%% 忐 %%%%%%%%%%
\subsection*{忐}\addcontentsline{loh}{figure}{忐}

\begin{Entry}{忐}{7}{⼼}
  \begin{Phonetics}{忐}{tan3}
    \definition{adj.}{nervoso; inquieto; agitado}
  \end{Phonetics}
\end{Entry}

%%%%%%%%%% 志 %%%%%%%%%%
\subsection*{志}\addcontentsline{loh}{figure}{志}

\begin{Entry}{志}{7}{⼼}
  \begin{Phonetics}{志}{zhi4}
    \definition{s.}{vontade; ideal; aspiração; ambição | anais; registros; transcrição | marca; sinal}
    \definition{v.}{ter em mente; lembrar | pesar; medir comprimento e quantidade}
  \end{Phonetics}
\end{Entry}

\begin{Entry}{志气}{7,4}{⼼,⽓}
  \begin{Phonetics}{志气}{zhi4qi4}[][HSK 7-9]
    \definition{s.}{ambição; aspiração; a determinação e a coragem para buscar a melhoria; a ambição de realizar algo}
  \synonymref{斗志}{dou4zhi4}
  \synonymref{骨气}{gu3qi4}
  \synonymref{理想}{li3xiang3}
  \synonymref{意向}{yi4xiang4}
  \synonymref{勇气}{yong3qi4}
  \synonymref{愿望}{yuan4wang4}
  \synonymref{志愿}{zhi4yuan4}
  \end{Phonetics}
\end{Entry}

\begin{Entry}{志愿}{7,14}{⼼,⽕}
  \begin{Phonetics}{志愿}{zhi4yuan4}[][HSK 3]
    \definition{s.}{desejo; ideal; aspiração; os ideais, desejos ou objetivos que se deseja realizar no coração}
    \definition{v.}{ser voluntário; ser proativo e disposto a realizar trabalhos sem remuneração ou com remuneração baixa, mas que possam ajudar outras pessoas}
  \end{Phonetics}
\end{Entry}

\begin{Entry}{志愿书}{7,14,4}{⼼,⽕,⼄}
  \begin{Phonetics}{志愿书}{zhi4yuan4shu1}
    \definition{s.}{formulário de inscrição; formulário de adesão; carta de intenções}
  \end{Phonetics}
\end{Entry}

\begin{Entry}{志愿者}{7,14,8}{⼼,⽕,⽼}
  \begin{Phonetics}{志愿者}{zhi4yuan4zhe3}[][HSK 3]
    \definition[名,位,个]{s.}{voluntário; pessoas que se voluntariam para prestar serviços em atividades sociais, grandes eventos esportivos, conferências, etc.}
  \end{Phonetics}
\end{Entry}

%%%%%%%%%% 忘 %%%%%%%%%%
\subsection*{忘}\addcontentsline{loh}{figure}{忘}

\begin{Entry}{忘}{7}{⼼}
  \begin{Phonetics}{忘}{wang4}[][HSK 1]
    \definition{v.}{esquecer | ignorar; negligenciar}
  \seealsoref{忘记}{wang4ji4}
  \seealsoref{忘却}{wang4que4}
  \synonymref{忘记}{wang4ji4}
  \antonymref{记}{ji4}
  \end{Phonetics}
\end{Entry}

\begin{Entry}{忘不了}{7,4,2}{⼼,⼀,⼅}
  \begin{Phonetics}{忘不了}{wang4bu5liao3}[][HSK 7-9]
    \definition{v.}{não poder esquecer}
  \end{Phonetics}
\end{Entry}

\begin{Entry}{忘本}{7,5}{⼼,⽊}
  \begin{Phonetics}{忘本}{wang4ben3}
    \definition{v.}{esquecer o sofrimento passado; esquecer de onde vem a felicidade | esquecer a própria origem de classe | esquecer a própria ancestralidade ou tradição}
  \synonymref{忘掉}{wang4/diao4}
  \synonymref{忘怀}{wang4huai2}
  \synonymref{忘记}{wang4ji4}
  \synonymref{忘却}{wang4que4}
  \synonymref{遗忘}{yi2wang4}
  \end{Phonetics}
\end{Entry}

\begin{Entry}{忘记}{7,5}{⼼,⾔}
  \begin{Phonetics}{忘记}{wang4ji4}[][HSK 1]
    \definition{v.}{esquecer | ignorar; negligenciar | sair da memória de alguém; não ser lembrado | descartar da mente; ignorar}
  \seealsoref{忘}{wang4}
  \seealsoref{忘却}{wang4que4}
  \synonymref{忘}{wang4}
  \synonymref{忘掉}{wang4/diao4}
  \synonymref{忘怀}{wang4huai2}
  \synonymref{忘却}{wang4que4}
  \synonymref{遗忘}{yi2wang4}
  \antonymref{惦记}{dian4ji4}
  \antonymref{怀念}{huai2nian4}
  \antonymref{记得}{ji4de5}
  \antonymref{纪念}{ji4nian4}
  \antonymref{记住}{ji4zhu5}
  \antonymref{牢记}{lao2ji4}
  \antonymref{铭记}{ming2ji4}
  \end{Phonetics}
\end{Entry}

\begin{Entry}{忘却}{7,7}{⼼,⼙}
  \begin{Phonetics}{忘却}{wang4que4}
    \definition{v.}{esquecer; cair no esquecimento}
  \seealsoref{忘掉}{wang4/diao4}
  \synonymref{忘掉}{wang4/diao4}
  \synonymref{忘怀}{wang4huai2}
  \synonymref{忘记}{wang4ji4}
  \synonymref{遗忘}{yi2wang4}
  \antonymref{怀念}{huai2nian4}
  \antonymref{记得}{ji4de5}
  \antonymref{纪念}{ji4nian4}
  \antonymref{记忆}{ji4yi4}
  \antonymref{牢记}{lao2ji4}
  \antonymref{缅怀}{mian3huai2}
  \antonymref{铭记}{ming2ji4}
  \antonymref{想念}{xiang3nian4}
  \end{Phonetics}
\end{Entry}

\begin{Entry}{忘怀}{7,7}{⼼,⼼}
  \begin{Phonetics}{忘怀}{wang4huai2}
    \definition{v.}{esquecer; apagar da mente}
  \synonymref{忘掉}{wang4/diao4}
  \synonymref{忘记}{wang4ji4}
  \synonymref{忘却}{wang4que4}
  \synonymref{遗忘}{yi2wang4}
  \antonymref{惦记}{dian4ji4}
  \antonymref{挂念}{gua4nian4}
  \antonymref{怀念}{huai2nian4}
  \antonymref{记得}{ji4de5}
  \antonymref{记住}{ji4zhu5}
  \antonymref{铭记}{ming2ji4}
  \antonymref{牵挂}{qian1gua4}
  \antonymref{想念}{xiang3nian4}
  \end{Phonetics}
\end{Entry}

\begin{Entry}{忘恩}{7,10}{⼼,⼼}
  \begin{Phonetics}{忘恩}{wang4'en1}
    \definition{v.}{ser ingrato}
  \synonymref{报仇}{bao4/chou2}
  \end{Phonetics}
\end{Entry}

\begin{Entry}{忘掉}{7,11}{⼼,⼿}
  \begin{Phonetics}{忘掉}{wang4/diao4}[][HSK 7-9]
    \definition{v.+compl.}{esquecer; deixar escapar da mente; não se lembrar}
  \synonymref{忘怀}{wang4huai2}
  \synonymref{忘记}{wang4ji4}
  \synonymref{忘却}{wang4que4}
  \synonymref{遗忘}{yi2wang4}
  \antonymref{记住}{ji4zhu5}
  \end{Phonetics}
\end{Entry}

\begin{Entry}{忘餐}{7,16}{⼼,⾷}
  \begin{Phonetics}{忘餐}{wang4can1}
    \definition{v.}{esquecer"-se das refeições}
  \end{Phonetics}
\end{Entry}

%%%%%%%%%% 忧 %%%%%%%%%%
\subsection*{忧}\addcontentsline{loh}{figure}{忧}

\begin{Entry}{忧}{7}{⼼}
  \begin{Phonetics}{忧}{you1}
    \definition{s.}{tristeza; ansiedade; preocupação; cuidado; coisas que causam tristeza}
    \definition{v.}{preocupar"-se; estar preocupado; estar ansioso; estar triste}
  \end{Phonetics}
\end{Entry}

\begin{Entry}{忧郁}{7,8}{⼼,⾢}
  \begin{Phonetics}{忧郁}{you1yu4}[][HSK 7-9]
    \definition{adj.}{abatido; melancólico; desanimado; triste}
    \definition{s.}{depressão | melancolia}
  \synonymref{担忧}{dan1you1}
  \synonymref{难过}{nan2guo4}
  \synonymref{忧愁}{you1chou2}
  \synonymref{忧虑}{you1lv4}
  \antonymref{高兴}{gao1xing4}
  \antonymref{开朗}{kai1lang3}
  \antonymref{愉快}{yu2kuai4}
  \end{Phonetics}
\end{Entry}

\begin{Entry}{忧虑}{7,10}{⼼,⾌}
  \begin{Phonetics}{忧虑}{you1lv4}[][HSK 7-9]
    \definition{adj.}{preocupado; ansioso; apreensivo}
    \definition[个,丝]{s.}{preocupação; apreensão; um estado de espírito ou sentimento de ansiedade}
  \synonymref{担心}{dan1/xin1}
  \synonymref{担忧}{dan1you1}
  \synonymref{顾虑}{gu4lv4}
  \synonymref{焦虑}{jiao1lv4}
  \synonymref{困扰}{kun4rao3}
  \synonymref{忧愁}{you1chou2}
  \synonymref{忧郁}{you1yu4}
  \synonymref{着急}{zhao2/ji2}
  \antonymref{安心}{an1xin1}
  \antonymref{放心}{fang4/xin1}
  \antonymref{快活}{kuai4huo5}
  \antonymref{期望}{qi1wang4}
  \antonymref{舒畅}{shu1chang4}
  \end{Phonetics}
\end{Entry}

\begin{Entry}{忧愁}{7,13}{⼼,⼼}
  \begin{Phonetics}{忧愁}{you1chou2}[][HSK 7-9]
    \definition{adj.}{triste; preocupado; deprimido; melancólico}
  \synonymref{担心}{dan1/xin1}
  \synonymref{担忧}{dan1you1}
  \synonymref{发愁}{fa1/chou2}
  \synonymref{烦闷}{fan2men4}
  \synonymref{烦恼}{fan2nao3}
  \synonymref{忧虑}{you1lv4}
  \synonymref{忧郁}{you1yu4}
  \antonymref{高兴}{gao1xing4}
  \antonymref{欢快}{huan1kuai4}
  \antonymref{欢乐}{huan1le4}
  \antonymref{开朗}{kai1lang3}
  \antonymref{快乐}{kuai4le4}
  \antonymref{快活}{kuai4huo5}
  \antonymref{痛快}{tong4kuai5}
  \antonymref{喜悦}{xi3yue4}
  \antonymref{欣喜}{xin1xi3}
  \end{Phonetics}
\end{Entry}

%%%%%%%%%% 快 %%%%%%%%%%
\subsection*{快}\addcontentsline{loh}{figure}{快}

\begin{Entry}{快}{7}{⼼}
  \begin{Phonetics}{快}{kuai4}[][HSK 1]
    \definition*{s.}{Sobrenome: Kuai}
    \definition{adj.}{rápido; veloz | apressado | perspicaz; ágil; inteligente; de mente rápida | (faca, espada, etc.) afiado | direto; franco; sem rodeios | satisfeito; feliz; gratificado | rápido; veloz; alta velocidade; tempo de execução curto | satisfeito; feliz; contente | engenhoso; ágil | afiado; facas, tesouras, machados e outros objetos afiados | sincero}
    \definition{adv.}{em breve; antes de muito tempo; estar prestes a | rapidamente}
    \definition{s.}{policial; polícia | Obsoleto: oficial encarregado de efetuar prisões}
  \seealsoref{快速}{kuai4su4}
  \seealsoref{迅速}{xun4su4}
  \synonymref{飞}{fei1}
  \synonymref{急}{ji2}
  \synonymref{疾}{ji2}
  \synonymref{利}{li4}
  \synonymref{灵}{ling2}
  \synonymref{勉}{mian3}
  \synonymref{敏}{min3}
  \synonymref{趣}{qu4}
  \synonymref{速}{su4}
  \antonymref{悲}{bei1}
  \antonymref{迟}{chi2}
  \antonymref{钝}{dun4}
  \antonymref{缓}{huan3}
  \antonymref{慢}{man4}
  \antonymref{伤}{shang1}
  \antonymref{滞}{zhi4}
  \antonymref{拙}{zhuo1}
  \end{Phonetics}
\end{Entry}

\begin{Entry}{快车}{7,4}{⼼,⾞}
  \begin{Phonetics}{快车}{kuai4che1}[][HSK 6]
    \definition{s.}{trem ou ônibus expresso; um trem ou ônibus com menos paradas e tempos de viagem mais curtos (usado principalmente para transporte de passageiros)}
  \antonymref{慢车}{man4che1}
  \end{Phonetics}
\end{Entry}

\begin{Entry}{快乐}{7,5}{⼼,⼃}
  \begin{Phonetics}{快乐}{kuai4le4}[][HSK 2]
    \definition{adj.}{feliz; alegre; animado; prazeiroso}
    \definition{s.}{felicidade | alegria}
  \seealsoref{快活}{kuai4huo5}
  \synonymref{高兴}{gao1xing4}
  \synonymref{开心}{kai1/xin1}
  \synonymref{快活}{kuai4huo5}
  \synonymref{痛快}{tong4kuai5}
  \synonymref{喜悦}{xi3yue4}
  \synonymref{兴奋}{xing1fen4}
  \synonymref{幸福}{xing4fu2}
  \synonymref{愉快}{yu2kuai4}
  \antonymref{悲哀}{bei1'ai1}
  \antonymref{悲伤}{bei1shang1}
  \antonymref{发愁}{fa1/chou2}
  \antonymref{难过}{nan2guo4}
  \antonymref{伤心}{shang1/xin1}
  \antonymref{痛苦}{tong4ku3}
  \antonymref{忧愁}{you1chou2}
  \antonymref{忧郁}{you1yu4}
  \end{Phonetics}
\end{Entry}

\begin{Entry}{快活}{7,9}{⼼,⽔}
  \begin{Phonetics}{快活}{kuai4huo5}[][HSK 5]
    \definition{adj.}{feliz; alegre; contente; animado}
  \end{Phonetics}
\end{Entry}

\begin{Entry}{快点儿}{7,9,2}{⼼,⽕,⼉}
  \begin{Phonetics}{快点儿}{kuai4dian3r5}[][HSK 2]
    \definition{v.}{apressar"-se}
  \end{Phonetics}
\end{Entry}

\begin{Entry}{快要}{7,9}{⼼,⾑}
  \begin{Phonetics}{快要}{kuai4yao4}[][HSK 2]
    \definition{adv.}{estar prestes a; estar indo para; estar à beira de; em breve; em pouco tempo; indica que a situação está prestes a ocorrer}
  \seealsoref{就要}{jiu4yao4}
  \synonymref{即将}{ji2jiang1}
  \synonymref{将近}{jiang1jin4}
  \synonymref{将要}{jiang1yao4}
  \synonymref{接近}{jie1jin4}
  \synonymref{就要}{jiu4yao4}
  \antonymref{已经}{yi3jing1}
  \end{Phonetics}
\end{Entry}

\begin{Entry}{快递}{7,10}{⼼,⾡}
  \begin{Phonetics}{快递}{kuai4di4}[][HSK 4]
    \definition[个,件,批]{s.}{correio rápido; entrega expressa; entrega rápida}
    \definition{v.}{entregar (serviço de entrega rápida por transportadoras especializadas)}
  \end{Phonetics}
\end{Entry}

\begin{Entry}{快速}{7,10}{⼼,⾡}
  \begin{Phonetics}{快速}{kuai4su4}[][HSK 3]
    \definition{adj.}{rápido; veloz; de alta velocidade; descreve o tempo curto gasto para caminhar, fazer algo, etc.}
  \end{Phonetics}
\end{Entry}

\begin{Entry}{快捷}{7,11}{⼼,⼿}
  \begin{Phonetics}{快捷}{kuai4jie2}[][HSK 7-9]
    \definition{adj.}{rápido; veloz (em relação à velocidade); ágil (em relação a pessoas)}
  \end{Phonetics}
\end{Entry}

\begin{Entry}{快餐}{7,16}{⼼,⾷}
  \begin{Phonetics}{快餐}{kuai4can1}[][HSK 2]
    \definition[份,顿]{s.}{pedido (comida) rápido; \emph{fast food}; refere"-se a refeições simples preparadas com antecedência e que podem ser servidas rapidamente}
  \synonymref{便饭}{bian4fan4}
  \end{Phonetics}
\end{Entry}

%%%%%%%%%% 怀 %%%%%%%%%%
\subsection*{怀}\addcontentsline{loh}{figure}{怀}

\begin{Entry}{怀}{7}{⼼}
  \begin{Phonetics}{怀}{huai2}
    \definition*{s.}{Sobrenome: Huai}
    \definition{s.}{seio; peito | mente}
    \definition{v.}{manter em mente; estimar; abrigar | sentir falta; pensar em; ansiar por | conceber (uma criança)}
  \end{Phonetics}
\end{Entry}

\begin{Entry}{怀孕}{7,5}{⼼,⼦}
  \begin{Phonetics}{怀孕}{huai2/yun4}[][HSK 7-9]
    \definition{v.+compl.}{estar (ficar) grávida}
  \end{Phonetics}
\end{Entry}

\begin{Entry}{怀旧}{7,5}{⼼,⽇}
  \begin{Phonetics}{怀旧}{huai2jiu4}[][HSK 7-9]
    \definition{s.}{lembrança afetuosa de tempos passados | nostalgia}
    \definition{v.}{sentir ou demonstrar nostalgia; lembrar de tempos passados ou de velhos conhecidos (geralmente com pensamentos gentis)}
  \end{Phonetics}
\end{Entry}

\begin{Entry}{怀里}{7,7}{⼼,⾥}
  \begin{Phonetics}{怀里}{huai2li5}[][HSK 7-9]
    \definition{s.}{seio; abraço}
  \end{Phonetics}
\end{Entry}

\begin{Entry}{怀念}{7,8}{⼼,⼼}
  \begin{Phonetics}{怀念}{huai2nian4}[][HSK 4]
    \definition{v.}{pensar em; valorizar a memória de}
  \end{Phonetics}
\end{Entry}

\begin{Entry}{怀抱}{7,8}{⼼,⼿}
  \begin{Phonetics}{怀抱}{huai2bao4}[][HSK 7-9]
    \definition{s.}{seio; peito | ambição; aspiração; intenção}
    \definition{v.}{abraçar; carregar nos braços; segurar nos braços | estimar; ter em mente}
  \end{Phonetics}
\end{Entry}

\begin{Entry}{怀着}{7,11}{⼼,⽬}
  \begin{Phonetics}{怀着}{huai2zhe5}[][HSK 7-9]
    \definition{v.}{nutrir; abrigar; ser preenchido com}
  \end{Phonetics}
\end{Entry}

\begin{Entry}{怀疑}{7,14}{⼼,⽦}
  \begin{Phonetics}{怀疑}{huai2yi2}[][HSK 4]
    \definition{v.}{duvidar; suspeitar | supor}
  \end{Phonetics}
\end{Entry}

%%%%%%%%%% 忠 %%%%%%%%%%
\subsection*{忠}\addcontentsline{loh}{figure}{忠}

\begin{Entry}{忠}{8}{⼼}
  \begin{Phonetics}{忠}{zhong1}
    \definition{adj.}{leal; fiel; devotado | honesto}
  \end{Phonetics}
\end{Entry}

\begin{Entry}{忠于}{8,3}{⼼,⼆}
  \begin{Phonetics}{忠于}{zhong1yu2}[][HSK 7-9]
    \definition{v.}{ser verdadeiro com; ser leal a; ser fiel a; ser devotado a}
  \synonymref{忠诚}{zhong1cheng2}
  \synonymref{忠心}{zhong1xin1}
  \synonymref{转载}{zhuan3zai3}
  \end{Phonetics}
\end{Entry}

\begin{Entry}{忠心}{8,4}{⼼,⼼}
  \begin{Phonetics}{忠心}{zhong1xin1}[][HSK 6]
    \definition{s.}{lealdade; devoção; fidelidade}
  \end{Phonetics}
\end{Entry}

\begin{Entry}{忠贞}{8,6}{⼼,⾙}
  \begin{Phonetics}{忠贞}{zhong1zhen1}[][HSK 7-9]
    \definition{adj.}{leal e inabalável}
    \definition{s.}{lealdade}
  \antonymref{变节}{bian4jie2}
  \end{Phonetics}
\end{Entry}

\begin{Entry}{忠实}{8,8}{⼼,⼧}
  \begin{Phonetics}{忠实}{zhong1shi2}[][HSK 7-9]
    \definition{adj.}{verdadeiro; leal; confiável; fiel; sincero | verdadeiro; verídico; autêntico; real}
  \synonymref{诚恳}{cheng2ken3}
  \synonymref{诚挚}{cheng2zhi4}
  \synonymref{诚实}{cheng2shi5}
  \synonymref{厚道}{hou4dao5}
  \synonymref{老实}{lao3shi5}
  \synonymref{忠诚}{zhong1cheng2}
  \antonymref{背叛}{bei4pan4}
  \antonymref{奸诈}{jian1zha4}
  \antonymref{虚伪}{xu1wei3}
  \end{Phonetics}
\end{Entry}

\begin{Entry}{忠诚}{8,8}{⼼,⾔}
  \begin{Phonetics}{忠诚}{zhong1cheng2}[][HSK 7-9]
    \definition{adj.}{leal; fiel; inabalável; (ao país, ao povo, à causa, aos amigos, etc.) fazer o melhor possível com a máxima sinceridade}
  \synonymref{忠实}{zhong1shi2}
  \antonymref{虚伪}{xu1wei3}
  \end{Phonetics}
\end{Entry}

%%%%%%%%%% 念 %%%%%%%%%%
\subsection*{念}\addcontentsline{loh}{figure}{念}

\begin{Entry}{念}{8}{⼼}
  \begin{Phonetics}{念}{nian4}[][HSK 3]
    \definition*{s.}{Sobrenome: Nian}
    \definition{num.}{vinte; 20; capitalização do número 廿}
    \definition{s.}{ideia; pensamento; pensamentos ou intenções internas}
    \definition{v.}{ler em voz alta | estudar; frequentar a escola | considerar; levar em conta | sentir falta; pensar em; pensar sobre; pensar frequentemente sobre}
  \seealsoref{廿}{nian4}
  \end{Phonetics}
\end{Entry}

\begin{Entry}{念书}{8,4}{⼼,⼄}
  \begin{Phonetics}{念书}{nian4/shu1}[][HSK 7-9]
    \definition{v.+compl.}{estudar; ir à escola; frequentar a escola | ler livros}
  \end{Phonetics}
\end{Entry}

\begin{Entry}{念头}{8,5}{⼼,⼤}
  \begin{Phonetics}{念头}{nian4tou5}[][HSK 7-9]
    \definition[个,种]{s.}{ideia; pensamento; intenção; plano}
  \end{Phonetics}
\end{Entry}

\begin{Entry}{念念不忘}{8,8,4,7}{⼼,⼼,⼀,⼼}
  \begin{Phonetics}{念念不忘}{nian4nian4-bu2wang4}[][HSK 7-9]
    \definition{expr.}{inesquecível; ter sempre em mente (expressão idiomática); nunca se esqueça}
  \end{Phonetics}
\end{Entry}

%%%%%%%%%% 忽 %%%%%%%%%%
\subsection*{忽}\addcontentsline{loh}{figure}{忽}

\begin{Entry}{忽}{8}{⼼}
  \begin{Phonetics}{忽}{hu1}
    \definition*{s.}{Sobrenome: Hu}
    \definition{adv.}{agora\dots, agora\dots | de repente; subitamente}[天气忽冷忽热。===O clima está frio em um minuto e quente no outro.]
    \definition{v.}{negligenciar; ignorar; não prestar atenção; não levar a sério}
  \end{Phonetics}
\end{Entry}

\begin{Entry}{忽视}{8,8}{⼼,⾒}
  \begin{Phonetics}{忽视}{hu1shi4}[][HSK 4]
    \definition{v.}{ignorar; negligenciar; menosprezar; desprezar; dar de ombros}
  \end{Phonetics}
\end{Entry}

\begin{Entry}{忽高忽低}{8,10,8,7}{⼼,⾼,⼼,⼈}
  \begin{Phonetics}{忽高忽低}{hu1gao1-hu1di1}[][HSK 7-9]
    \definition{expr.}{``Altos e baixos.''; ora alto, ora baixo}
  \end{Phonetics}
\end{Entry}

\begin{Entry}{忽悠}{8,11}{⼼,⼼}
  \begin{Phonetics}{忽悠}{hu1you5}[][HSK 7-9]
    \definition{v.}{balançar; cintilar; sacudir | enganar; enganar alguém}
  \end{Phonetics}
\end{Entry}

\begin{Entry}{忽略}{8,11}{⼼,⽥}
  \begin{Phonetics}{忽略}{hu1lve4}[][HSK 6]
    \definition{v.}{negligenciar; ignorar; não perceber}
  \end{Phonetics}
\end{Entry}

\begin{Entry}{忽然}{8,12}{⼼,⽕}
  \begin{Phonetics}{忽然}{hu1ran2}[][HSK 2]
    \definition{adv.}{repentinamente; de repente; sem aviso prévio; significa que algo aconteceu de forma rápida e inesperada}
  \seealsoref{突然}{tu1ran2}
  \synonymref{猛然}{meng3ran2}
  \synonymref{突然}{tu1ran2}
  \synonymref{骤然}{zhou4ran2}
  \antonymref{渐渐}{jian4jian4}
  \antonymref{慢慢}{man4man4}
  \antonymref{逐步}{zhu2bu4}
  \antonymref{逐渐}{zhu2jian4}
  \end{Phonetics}
\end{Entry}

%%%%%%%%%% 态 %%%%%%%%%%
\subsection*{态}\addcontentsline{loh}{figure}{态}

\begin{Entry}{态}{8}{⼼}
  \begin{Phonetics}{态}{tai4}
    \definition{s.}{forma; aparência; condição | (física) estado | (linguística) voz}[气态===estado gasoso | 被动态===voz passiva]
  \end{Phonetics}
\end{Entry}

\begin{Entry}{态度}{8,9}{⼼,⼴}
  \begin{Phonetics}{态度}{tai4du5}[][HSK 2]
    \definition[种,个]{s.}{maneira; comportamento; atitude; comportamento e expressão facial das pessoas | atitude; abordagem; opinião sobre o assunto e medidas tomadas}
  \synonymref{看法}{kan4fa5}
  \synonymref{立场}{li4chang3}
  \synonymref{神情}{shen2qing2}
  \synonymref{神态}{shen2tai4}
  \synonymref{心态}{xin1tai4}
  \synonymref{姿势}{zi1shi4}
  \synonymref{作风}{zuo4feng1}
  \end{Phonetics}
\end{Entry}

%%%%%%%%%% 怕 %%%%%%%%%%
\subsection*{怕}\addcontentsline{loh}{figure}{怕}

\begin{Entry}{怕}{8}{⼼}
  \begin{Phonetics}{怕}{pa4}[][HSK 2]
    \definition{adv.}{(expressando suposição, julgamento, estimativa, etc.) talvez; suponho; receio (que)}
    \definition{adv.}{por medo; talvez; suponho}
    \definition{v.}{temer; ter medo; recear; sentir medo, ficar nervoso | estar preocupado com; estar preocupado por (ou sobre); ter medo de que algo possa acontecer | ser afetado por; não conseguir suportar; não aguentar mais}
  \seealsoref{害怕}{hai4pa4}
  \seealsoref{恐怕}{kong3pa4}
  \synonymref{害怕}{hai4pa4}
  \synonymref{恐怕}{kong3pa4}
  \synonymref{畏}{wei4}
  \end{Phonetics}
\end{Entry}

%%%%%%%%%% 怛 %%%%%%%%%%
\subsection*{怛}\addcontentsline{loh}{figure}{怛}

\begin{Entry}{怛}{8}{⼼}
  \begin{Phonetics}{怛}{da2}
    \definition{adj.}{Arcaico: aflito; angustiado | Arcaico: aterrorizado | Arcaico: alarmado; entristecido | Arcaico: chocado}
  \end{Phonetics}
\end{Entry}

%%%%%%%%%% 怜 %%%%%%%%%%
\subsection*{怜}\addcontentsline{loh}{figure}{怜}

\begin{Entry}{怜}{8}{⼼}
  \begin{Phonetics}{怜}{lian2}
    \definition{v.}{simpatizar com; ter pena; sentir compaixão}
  \end{Phonetics}
\end{Entry}

\begin{Entry}{怜惜}{8,11}{⼼,⼼}
  \begin{Phonetics}{怜惜}{lian2xi1}[][HSK 7-9]
    \definition{v.}{ter pena de; sentir compaixão por | sentir ternura e proteção em relação a}
  \end{Phonetics}
\end{Entry}

%%%%%%%%%% 怡 %%%%%%%%%%
\subsection*{怡}\addcontentsline{loh}{figure}{怡}

\begin{Entry}{怡}{8}{⼼}
  \begin{Phonetics}{怡}{yi2}
    \definition*{s.}{Sobrenome: Yi}
    \definition{adj.}{Literário: feliz; alegre; animado}
  \end{Phonetics}
\end{Entry}

\begin{Entry}{怡然自得}{8,12,6,11}{⼼,⽕,⾃,⼻}
  \begin{Phonetics}{怡然自得}{yi2ran2-zi4de2}[][HSK 7-9]
    \definition{expr.}{feliz e contente; ser feliz e estar satisfeito consigo mesmo; sentir uma sensação de felicidade plena}
  \end{Phonetics}
\end{Entry}

%%%%%%%%%% 性 %%%%%%%%%%
\subsection*{性}\addcontentsline{loh}{figure}{性}

\begin{Entry}{性}{8}{⼼}
  \begin{Phonetics}{性}{xing4}[][HSK 3]
    \definition[个]{s.}{natureza; caráter; personalidade | propriedade; qualidade; natureza e características das coisas | sexo; gênero | sexualidade; relacionado com a reprodução e a sexualidade | caráter; temperamento}
    \definition{suf.}{indica uma determinada propriedade ou característica de algo; segue um substantivo, verbo ou adjetivo, formando um substantivo abstrato ou um adjetivo que expressa uma propriedade}
  \end{Phonetics}
\end{Entry}

\begin{Entry}{性生活}{8,5,9}{⼼,⽣,⽔}
  \begin{Phonetics}{性生活}{xing4sheng1huo2}
    \definition{s.}{vida sexual}
  \end{Phonetics}
\end{Entry}

\begin{Entry}{性价比}{8,6,4}{⼼,⼈,⽐}
  \begin{Phonetics}{性价比}{xing4jia4bi3}[][HSK 7-9]
    \definition{s.}{relação custo"-benefício; a relação entre o valor de desempenho de um produto e seu valor de preço é uma medida quantitativa que reflete a acessibilidade do produto | desempenho de custo}
  \end{Phonetics}
\end{Entry}

\begin{Entry}{性别}{8,7}{⼼,⼑}
  \begin{Phonetics}{性别}{xing4bie2}[][HSK 3]
    \definition[种]{s.}{sexo; gênero}
  \end{Phonetics}
\end{Entry}

\begin{Entry}{性命}{8,8}{⼼,⼝}
  \begin{Phonetics}{性命}{xing4ming4}[][HSK 7-9]
    \definition[条]{s.}{vida (de um homem ou animal)}
  \seealsoref{生命}{sheng1ming4}
  \synonymref{生命}{sheng1ming4}
  \end{Phonetics}
\end{Entry}

\begin{Entry}{性质}{8,8}{⼼,⾙}
  \begin{Phonetics}{性质}{xing4zhi4}[][HSK 4]
    \definition[个,种,类]{s.}{natureza; qualidade; caráter; propriedade; propriedade fundamental que distingue uma coisa de outra}
  \end{Phonetics}
\end{Entry}

\begin{Entry}{性侵}{8,9}{⼼,⼈}
  \begin{Phonetics}{性侵}{xing4qin1}
    \definition{s.}{agressão sexual}
    \definition{v.}{agredir sexualmente}
  \end{Phonetics}
\end{Entry}

\begin{Entry}{性格}{8,10}{⼼,⽊}
  \begin{Phonetics}{性格}{xing4ge2}[][HSK 3]
    \definition[种,个]{s.}{caráter; temperamento; as características psicológicas manifestadas na atitude e no comportamento em relação às pessoas e às coisas}
  \end{Phonetics}
\end{Entry}

\begin{Entry}{性能}{8,10}{⼼,⾁}
  \begin{Phonetics}{性能}{xing4neng2}[][HSK 5]
    \definition{s.}{natureza; propriedade; desempenho; função (de uma máquina, etc.); grau de conformidade dos produtos mecânicos ou outros produtos industriais com os requisitos de projeto}
  \end{Phonetics}
\end{Entry}

\begin{Entry}{性情}{8,11}{⼼,⼼}
  \begin{Phonetics}{性情}{xing4qing2}[][HSK 7-9]
    \definition[种]{s.}{disposição; temperamento; personalidade; disposição e caráter | sentimento; emoções; natureza interior; caráter emocional; pensamentos e sentimentos}
  \seealsoref{性格}{xing4ge2}
  \synonymref{本性}{ben3xing4}
  \synonymref{个性}{ge4xing4}
  \synonymref{脾气}{pi2qi5}
  \synonymref{特性}{te4xing4}
  \synonymref{天性}{tian1xing4}
  \synonymref{性格}{xing4ge2}
  \end{Phonetics}
\end{Entry}

%%%%%%%%%% 怪 %%%%%%%%%%
\subsection*{怪}\addcontentsline{loh}{figure}{怪}

\begin{Entry}{怪}{8}{⼼}
  \begin{Phonetics}{怪}{guai4}[][HSK 4,5]
    \definition*{s.}{Sobrenome: Guai}
    \definition{adj.}{estranho; esquisito; desconcertante | peculiar; excêntrico; pitoresco; monstruoso; anormal; incomum}
    \definition{adv.}{bastante; muito}
    \definition{s.}{monstro; demônio | diabo; ser maligno}
    \definition{v.}{culpar | achar algo estranho; maravilhar"-se com; ficar surpreso | repreender; culpar; reclamar}
  \end{Phonetics}
\end{Entry}

\begin{Entry}{怪不得}{8,4,11}{⼼,⼀,⼻}
  \begin{Phonetics}{怪不得}{guai4bu5de5}[][HSK 7-9]
    \definition{adv.}{não é de admirar; então é por isso; isso explica por que; isso significa que você entende o motivo e não acha mais uma situação estranha}
    \definition{v.}{não culpar; não acusar; não poder culpar, não se ofender}[你做错了，怪不得别人。===Você cometeu um erro, então não culpe os outros.]
  \end{Phonetics}
\end{Entry}

\begin{Entry}{怪异}{8,6}{⼼,⼶}
  \begin{Phonetics}{怪异}{guai4yi4}[][HSK 7-9]
    \definition{adj.}{monstruoso; estranho; incomum}
    \definition{s.}{fenômeno estranho; presságio; prodígio | monstruosidade}
  \end{Phonetics}
\end{Entry}

\begin{Entry}{怪物}{8,8}{⼼,⽜}
  \begin{Phonetics}{怪物}{guai4wu5}[][HSK 7-9]
    \definition{s.}{monstro; aberração; coisas imaginárias que parecem estranhas, mas têm habilidades especiais | pessoa excêntrica; pássaro estranho; uma pessoa com temperamento excêntrico}
  \end{Phonetics}
\end{Entry}

\begin{Entry}{怪兽}{8,11}{⼼,⼋}
  \begin{Phonetics}{怪兽}{guai4shou4}
    \definition{s.}{monstro; criaturas míticas com formas bizarras}
  \end{Phonetics}
\end{Entry}

\begin{Entry}{怪癖}{8,18}{⼼,⽧}
  \begin{Phonetics}{怪癖}{guai4pi3}
    \definition{s.}{excentricidade; manias | peculiaridade | hobby estranho}
  \synonymref{古怪}{gu3guai4}
  \end{Phonetics}
\end{Entry}

%%%%%%%%%% 怎 %%%%%%%%%%
\subsection*{怎}\addcontentsline{loh}{figure}{怎}

\begin{Entry}{怎}{9}{⼼}
  \begin{Phonetics}{怎}{zen3}
    \definition{pron.}{por que?; como?; expressa uma pergunta, equivalente a 怎么}
  \seealsoref{怎么}{zen3me5}
  \end{Phonetics}
\end{Entry}

\begin{Entry}{怎么}{9,3}{⼼,⼃}
  \begin{Phonetics}{怎么}{zen3me5}[][HSK 1]
    \definition{pron.}{como?; o quê?; perguntas sobre natureza, situação, método, motivo, etc. | de qualquer maneira; não importa como; de uma certa maneira; referência geral à natureza, condição ou modo | que? (usado sozinho no início de uma frase para expressar surpresa) | usado após 不 e 没, indica um grau baixo e é uma forma mais educada de se expressar | usado em perguntas retóricas}
  \seealsoref{不}{bu4}
  \seealsoref{没}{mei2}
  \seealsoref{怎么样}{zen3me5yang4}
  \seealsoref{怎样}{zen3yang4}
  \synonymref{如何}{ru2he2}
  \synonymref{怎么样}{zen3me5yang4}
  \synonymref{怎样}{zen3yang4}
  \end{Phonetics}
\end{Entry}

\begin{Entry}{怎么了}{9,3,2}{⼼,⼃,⼅}
  \begin{Phonetics}{怎么了}{zen3me5 le5}
    \definition{expr.}{``E aí?''; ``O que aconteceu?''; ``O que está rolando?''; informe"-se sobre a situação}
  \end{Phonetics}
\end{Entry}

\begin{Entry}{怎么办}{9,3,4}{⼼,⼃,⼒}
  \begin{Phonetics}{怎么办}{zen3me5ban4}[][HSK 2]
    \definition{adv.}{o que fazer?; o que deve ser feito?}
  \end{Phonetics}
\end{Entry}

\begin{Entry}{怎么回事}{9,3,6,8}{⼼,⼃,⼞,⼅}
  \begin{Phonetics}{怎么回事}{zen3me5hui2shi4}
    \definition{expr.}{``Como isso é possível?''; usado para indagar sobre a causa ou situação de algo | ``Como isso aconteceu?''; usado para perguntar sobre o progresso ou o estado de algo}
  \end{Phonetics}
\end{Entry}

\begin{Entry}{怎么样}{9,3,10}{⼼,⼃,⽊}
  \begin{Phonetics}{怎么样}{zen3me5yang4}[][HSK 2]
    \definition{adv.}{como?; o que?; como é?; como estão as coisas?; o que você acha?; pergunte sobre o método, natureza, situação, opinião, etc. | substitui uma ação ou situação não dita (usado apenas na forma negativa, mais eufemístico do que uma declaração direta); indaga sobre a natureza, condição, método, razão, etc.}
  \seealsoref{什么样}{shen2me5yang4}
  \seealsoref{怎么}{zen3me5}
  \synonymref{如何}{ru2he2}
  \synonymref{什么样}{shen2me5yang4}
  \synonymref{怎么}{zen3me5}
  \synonymref{怎样}{zen3yang4}
  \end{Phonetics}
\end{Entry}

\begin{Entry}{怎么得了}{9,3,11,2}{⼼,⼃,⼻,⼅}
  \begin{Phonetics}{怎么得了}{zen3me5de2liao3}
    \definition{expr.}{``O que está acontecendo?'';  ``Que coisa terrível seria!";  ``Como isso é possível?''; ``Que bagunça horrível!'';  ``O que deve ser feito?''}
  \end{Phonetics}
\end{Entry}

\begin{Entry}{怎么搞的}{9,3,13,8}{⼼,⼃,⼿,⽩}
  \begin{Phonetics}{怎么搞的}{zen3me5 gao3 de5}
    \definition{expr.}{``Como isso aconteceu?'' | ``O que deu errado?'' | ``E aí?''; ``Qual é?'' | ``O que há de errado?''}
  \end{Phonetics}
\end{Entry}

\begin{Entry}{怎样}{9,10}{⼼,⽊}
  \begin{Phonetics}{怎样}{zen3yang4}[][HSK 2]
    \definition{pron.}{como?; o que?; indagar sobre a natureza, condição ou método, etc. | como?; indica uma referência virtual | de uma certa maneira; de qualquer maneira; não importa como; indica qualquer | como?; usado como predicado, objeto ou complemento para indagar sobre uma situação}
  \seealsoref{怎么}{zen3me5}
  \synonymref{如何}{ru2he2}
  \synonymref{怎么}{zen3me5}
  \end{Phonetics}
\end{Entry}

%%%%%%%%%% 怒 %%%%%%%%%%
\subsection*{怒}\addcontentsline{loh}{figure}{怒}

\begin{Entry}{怒}{9}{⼼}
  \begin{Phonetics}{怒}{nu4}
    \definition{adj.}{zangado; furioso | feroz; forte; descreve um forte impulso}
    \definition{adv.}{com força; vigorosamente; dinamicamente | com raiva}
    \definition{s.}{raiva; fúria}
    \definition{v.}{enfurecer"-se; ficar com raiva}
  \end{Phonetics}
\end{Entry}

\begin{Entry}{怒放}{9,8}{⼼,⽅}
  \begin{Phonetics}{怒放}{nu4fang4}
    \definition{v.}{florescer em plena floração}
  \end{Phonetics}
\end{Entry}

\begin{Entry}{怒骂}{9,9}{⼼,⾺}
  \begin{Phonetics}{怒骂}{nu4ma4}
    \definition{v.}{xingar em fúria | praguejar furiosamente; gritar insultos raivosos}
  \synonymref{痛骂}{tong4ma4}
  \antonymref{嬉笑}{xi1xiao4}
  \end{Phonetics}
\end{Entry}

%%%%%%%%%% 思 %%%%%%%%%%
\subsection*{思}\addcontentsline{loh}{figure}{思}

\begin{Entry}{思}{9}{⼼}
  \begin{Phonetics}{思}{si1}
    \definition*{s.}{Sobrenome: Si}
    \definition{s.}{pensamento; ideias | pensamentos; emoções; humor}
    \definition{v.}{pensar; considerar; deliberar | pensar em; ansiar por}
  \end{Phonetics}
\end{Entry}

\begin{Entry}{思考}{9,6}{⼼,⽼}
  \begin{Phonetics}{思考}{si1kao3}[][HSK 4]
    \definition{v.}{pensar; ponderar; considerar; deliberar; envolver"-se em atividades de pensamento, como análise, síntese, julgamento, raciocínio e generalização}
  \end{Phonetics}
\end{Entry}

\begin{Entry}{思念}{9,8}{⼼,⼼}
  \begin{Phonetics}{思念}{si1nian4}[][HSK 7-9]
    \definition{v.}{sentir falta; pensar em; ansiar por; sentir saudades}
  \synonymref{怀念}{huai2nian4}
  \end{Phonetics}
\end{Entry}

\begin{Entry}{思前想后}{9,9,13,6}{⼼,⼑,⼼,⼝}
  \begin{Phonetics}{思前想后}{si1qian2-xiang3hou4}[][HSK 7-9]
    \definition{expr.}{``Depois de refletir bastante sobre o assunto.''; refletir repetidamente; ponderar sobre a mesma coisa; considerar a causa passada e o efeito futuro; refletir sobre o passado e o futuro; refletir sobre as razões e a conexão}
  \end{Phonetics}
\end{Entry}

\begin{Entry}{思索}{9,10}{⼼,⽷}
  \begin{Phonetics}{思索}{si1suo3}[][HSK 7-9]
    \definition{v.}{ponderar; considerar; especular; deliberar; pensar profundamente}
  \synonymref{考虑}{kao3lv4}
  \synonymref{思考}{si1kao3}
  \synonymref{思念}{si1nian4}
  \synonymref{推敲}{tui1qiao1}
  \synonymref{研究}{yan2jiu1}
  \antonymref{行动}{xing2dong4}
  \end{Phonetics}
\end{Entry}

\begin{Entry}{思维}{9,11}{⼼,⽷}
  \begin{Phonetics}{思维}{si1wei2}[][HSK 5]
    \definition[种]{s.}{pensamento; reflexão; organizar e transformar os materiais obtidos através do conhecimento sensorial para formar conceitos, julgamentos e raciocínios}
    \definition{v.}{pensar}
  \end{Phonetics}
\end{Entry}

\begin{Entry}{思想}{9,13}{⼼,⼼}
  \begin{Phonetics}{思想}{si1xiang3}[][HSK 3]
    \definition[个,种]{s.}{reflexão; pensamento; ideologia; a existência objetiva é refletida na consciência das pessoas por meio de atividades de pensamento, que pertencem à cognição racional | ideia; pensamento}
  \end{Phonetics}
\end{Entry}

\begin{Entry}{思路}{9,13}{⼼,⾜}
  \begin{Phonetics}{思路}{si1lu4}[][HSK 7-9]
    \definition[种]{s.}{ideias; pensamento; mentalidade; linha de raciocínio}
  \synonymref{灵感}{ling2gan3}
  \synonymref{线索}{xian4suo3}
  \end{Phonetics}
\end{Entry}

%%%%%%%%%% 怠 %%%%%%%%%%
\subsection*{怠}\addcontentsline{loh}{figure}{怠}

\begin{Entry}{怠}{9}{⼼}
  \begin{Phonetics}{怠}{dai4}
    \definition{adj.}{ocioso; relaxado; negligente | preguiçoso; indolente}
    \definition{v.}{ficar ocioso; negligenciar; folgar}
  \end{Phonetics}
\end{Entry}

\begin{Entry}{怠工}{9,3}{⼼,⼯}
  \begin{Phonetics}{怠工}{dai4/gong1}[][HSK 7-9]
    \definition{v.+compl.}{ir devagar (como uma forma de greve); ir devagar | afrouxar no trabalho}
  \end{Phonetics}
\end{Entry}

\begin{Entry}{怠慢}{9,14}{⼼,⼼}
  \begin{Phonetics}{怠慢}{dai4man4}[][HSK 7-9]
    \definition{v.}{ignorar; tratar com indiferença | deixar de dar a devida atenção para alguém}
  \end{Phonetics}
\end{Entry}

%%%%%%%%%% 急 %%%%%%%%%%
\subsection*{急}\addcontentsline{loh}{figure}{急}

\begin{Entry}{急}{9}{⼼}
  \begin{Phonetics}{急}{ji2}[][HSK 2]
    \definition{adj.}{impaciente; ansioso | irritado; aborrecido; incomodado | rápido e intenso; veloz | urgente; premente}
    \definition{s.}{urgência; emergência; assunto urgente e grave}
    \definition{v.}{preocupar; deixar ansioso | estar ansioso para ajudar; tratar os problemas dos outros como se fossem urgentes e ajudar a resolvê"-los imediatamente}
  \seealsoref{着急}{zhao2/ji2}
  \synonymref{匆}{cong1}
  \synonymref{快}{kuai4}
  \synonymref{着急}{zhao2/ji2}
  \antonymref{缓}{huan3}
  \antonymref{慢}{man4}
  \end{Phonetics}
\end{Entry}

\begin{Entry}{急于}{9,3}{⼼,⼆}
  \begin{Phonetics}{急于}{ji2yu2}[][HSK 7-9]
    \definition{v.}{estar (ser) ansioso; estar (ser) impaciente; estar ansioso para}
  \end{Phonetics}
\end{Entry}

\begin{Entry}{急忙}{9,6}{⼼,⼼}
  \begin{Phonetics}{急忙}{ji2mang2}[][HSK 4]
    \definition{adv.}{apressadamente; com pressa}
  \end{Phonetics}
\end{Entry}

\begin{Entry}{急诊}{9,7}{⼼,⾔}
  \begin{Phonetics}{急诊}{ji2zhen3}[][HSK 7-9]
    \definition{s.}{pronto"-socorro; emergência; tratamento de emergência; uma clínica ambulatorial especial em um hospital para pessoas com doenças agudas}
  \end{Phonetics}
\end{Entry}

\begin{Entry}{急性}{9,8}{⼼,⼼}
  \begin{Phonetics}{急性}{ji2xing4}[][HSK 7-9]
    \definition{adj.}{aguda}
    \definition{s.}{pessoa impetuosa; cabeça quente}
  \seealsoref{急性儿}{ji2xing4r5}
  \antonymref{慢性}{man4xing4}
  \end{Phonetics}
\end{Entry}

\begin{Entry}{急性儿}{9,8,2}{⼼,⼼,⼉}
  \begin{Phonetics}{急性儿}{ji2xing4r5}
    \definition{adj.}{impetuoso; temperamental; de temperamento explosivo; de disposição impaciente}
  \end{Phonetics}
\end{Entry}

\begin{Entry}{急转弯}{9,8,9}{⼼,⾞,⼸}
  \begin{Phonetics}{急转弯}{ji2zhuan3wan1}[][HSK 7-9]
    \definition{s.}{curva fechada; curva acetuada; cotovelo}
    \definition{v.}{Coloquial: (uma atitude, política, etc.) fazer uma mudança repentina | fazer uma curva repentina | fazer uma mudança radical}
  \seealsoref{急转弯儿}{ji2zhuan3wan1r5}
  \end{Phonetics}
\end{Entry}

\begin{Entry}{急转弯儿}{9,8,9,2}{⼼,⾞,⼸,⼉}
  \begin{Phonetics}{急转弯儿}{ji2zhuan3wan1r5}
    \definition{s.}{curva fechada}
  \end{Phonetics}
\end{Entry}

\begin{Entry}{急迫}{9,8}{⼼,⾡}
  \begin{Phonetics}{急迫}{ji2po4}[][HSK 7-9]
    \definition{adj.}{urgente; premente; imperativo}
  \end{Phonetics}
\end{Entry}

\begin{Entry}{急剧}{9,10}{⼼,⼑}
  \begin{Phonetics}{急剧}{ji2ju4}[][HSK 7-9]
    \definition{adj.}{rápido; agudo; repentino}
    \definition{adv.}{Formal: rapidamente (geralmente mudando em uma direção ruim ou potencialmente levando a resultados ruins)}
  \end{Phonetics}
\end{Entry}

\begin{Entry}{急救}{9,11}{⼼,⽁}
  \begin{Phonetics}{急救}{ji2jiu4}[][HSK 6]
    \definition{s.}{primeiros socorros; tratamento médico de emergência (para pessoas gravemente doentes ou gravemente feridas)}
    \definition{v.}{prestar primeiros socorros; dar tratamento de emergência}
  \end{Phonetics}
\end{Entry}

\begin{Entry}{急需}{9,14}{⼼,⾬}
  \begin{Phonetics}{急需}{ji2xu1}[][HSK 7-9]
    \definition{v.}{estar em extrema necessidade de}
  \end{Phonetics}
\end{Entry}

%%%%%%%%%% 怨 %%%%%%%%%%
\subsection*{怨}\addcontentsline{loh}{figure}{怨}

\begin{Entry}{怨}{9}{⼼}
  \begin{Phonetics}{怨}{yuan4}[][HSK 5]
    \definition{s.}{ressentimento; inimizade; rancor}
    \definition{v.}{culpar; reclamar}
  \end{Phonetics}
\end{Entry}

\begin{Entry}{怨气}{9,4}{⼼,⽓}
  \begin{Phonetics}{怨气}{yuan4qi4}[][HSK 7-9]
    \definition{s.}{queixa; reclamação; ressentimento; uma expressão ou emoção de ressentimento}
  \antonymref{和气}{he2qi5}
  \end{Phonetics}
\end{Entry}

\begin{Entry}{怨言}{9,7}{⼼,⾔}
  \begin{Phonetics}{怨言}{yuan4yan2}[][HSK 7-9]
    \definition{s.}{reclamação; resmungo}
  \synonymref{抱怨}{bao4yuan5}
  \antonymref{赞美}{zan4mei3}
  \end{Phonetics}
\end{Entry}

\begin{Entry}{怨恨}{9,9}{⼼,⼼}
  \begin{Phonetics}{怨恨}{yuan4hen4}[][HSK 7-9]
    \definition{s.}{ressentimento; forte insatisfação ou ódio; intensa insatisfação ou ódio}
    \definition{v.}{sentir ressentimento; guardar rancor de alguém; forte insatisfação ou ódio em relação a uma pessoa ou coisa}
  \synonymref{抱怨}{bao4yuan5}
  \synonymref{仇恨}{chou2hen4}
  \synonymref{后悔}{hou4hui3}
  \synonymref{悔恨}{hui3hen4}
  \synonymref{埋怨}{man2yuan4}
  \antonymref{恩惠}{en1hui4}
  \antonymref{感激}{gan3ji1}
  \antonymref{宽恕}{kuan1shu4}
  \end{Phonetics}
\end{Entry}

%%%%%%%%%% 怹 %%%%%%%%%%
\subsection*{怹}\addcontentsline{loh}{figure}{怹}

\begin{Entry}{怹}{9}{⼼}
  \begin{Phonetics}{怹}{tan1}
    \definition{pron.}{ele, ela (cortês)}
  \synonymref{他}{ta1}
  \end{Phonetics}
\end{Entry}

%%%%%%%%%% 总 %%%%%%%%%%
\subsection*{总}\addcontentsline{loh}{figure}{总}

\begin{Entry}{总}{9}{⼼}
  \begin{Phonetics}{总}{zong3}[][HSK 3]
    \definition{adj.}{total; geral; global | responsável (liderança)}
    \definition{adv.}{sempre; invariavelmente | de qualquer forma; afinal; eventualmente; mais cedo ou mais tarde; no fim das contas | certamente; provavelmente; com certeza; expressa estimativa; suposição; equivalente a 大概}
    \definition{v.}{reunir; resumir; juntar; compilar}
  \seealsoref{大概}{da4gai4}
  \end{Phonetics}
\end{Entry}

\begin{Entry}{总之}{9,3}{⼼,⼂}
  \begin{Phonetics}{总之}{zong3zhi1}[][HSK 4]
    \definition{conj.}{em uma palavra; em suma; em resumo; indica que a declaração seguinte é uma declaração geral}
  \end{Phonetics}
\end{Entry}

\begin{Entry}{总计}{9,4}{⼼,⾔}
  \begin{Phonetics}{总计}{zong3ji4}[][HSK 7-9]
    \definition{s.}{total}
    \definition{v.}{totalizar; equivaler a; somar até}
  \synonymref{共计}{gong4ji4}
  \synonymref{合计}{he2ji5}
  \synonymref{累计}{lei3ji4}
  \synonymref{全部}{quan2bu4}
  \synonymref{统计}{tong3ji4}
  \synonymref{一共}{yi2gong4}
  \synonymref{总共}{zong3gong4}
  \antonymref{单独}{dan1du2}
  \end{Phonetics}
\end{Entry}

\begin{Entry}{总长}{9,4}{⼼,⾧}
  \begin{Phonetics}{总长}{zong3chang2}
    \definition{s.}{comprimento total}
  \end{Phonetics}
  \begin{Phonetics}{总长}{zong3zhang3}
    \definition{s.}{nome usado para ministros de gabinete entre 1912 e 1927, substituído por 部长}
  \seealsoref{部长}{bu4zhang3}
  \end{Phonetics}
\end{Entry}

\begin{Entry}{总务}{9,5}{⼼,⼒}
  \begin{Phonetics}{总务}{zong3wu4}
    \definition{s.}{assuntos gerais | pessoa encarregada dos assuntos gerais; pessoa responsável pelas tarefas administrativas | serviços gerais; tarefas administrativas em agências governamentais e escolas}
  \end{Phonetics}
\end{Entry}

\begin{Entry}{总台}{9,5}{⼼,⼝}
  \begin{Phonetics}{总台}{zong3tai2}
    \definition{s.}{balcão de recepção; peitoril da janela}
  \end{Phonetics}
\end{Entry}

\begin{Entry}{总价}{9,6}{⼼,⼈}
  \begin{Phonetics}{总价}{zong3jia4}
    \definition{s.}{preço total; valor total, o valor que constitui o número ou a quantidade inteira}
  \end{Phonetics}
\end{Entry}

\begin{Entry}{总共}{9,6}{⼼,⼋}
  \begin{Phonetics}{总共}{zong3gong4}[][HSK 4]
    \definition{adv.}{em tudo; em todos; no total; completamente; totalmente; em conjunto}
  \end{Phonetics}
\end{Entry}

\begin{Entry}{总而言之}{9,6,7,3}{⼼,⽽,⾔,⼂}
  \begin{Phonetics}{总而言之}{zong3'er2yan2zhi1}[][HSK 7-9]
    \definition{expr.}{resumindo; em poucas palavras; em suma; para encurtar uma longa história}
  \synonymref{综上所述}{zong1shang4-suo3shu4}
  \end{Phonetics}
\end{Entry}

\begin{Entry}{总体}{9,7}{⼼,⼈}
  \begin{Phonetics}{总体}{zong3ti3}[][HSK 5]
    \definition{s.}{total; geral; conjunto; totalidade; massa; população; o todo formado pela união de vários indivíduos; a totalidade das coisas}
  \end{Phonetics}
\end{Entry}

\begin{Entry}{总的来说}{9,8,7,9}{⼼,⽩,⽊,⾔}
  \begin{Phonetics}{总的来说}{zong3de5lai2shuo1}[][HSK 7-9]
    \definition{adv.}{em geral; de modo geral; como um grupo; em suma; resumindo; para concluir}
  \end{Phonetics}
\end{Entry}

\begin{Entry}{总线}{9,8}{⼼,⽷}
  \begin{Phonetics}{总线}{zong3xian4}
    \definition{s.}{barramento (computador); \emph{computer bus}}
  \end{Phonetics}
\end{Entry}

\begin{Entry}{总经理}{9,8,11}{⼼,⽷,⽟}
  \begin{Phonetics}{总经理}{zong3jing1li3}[][HSK 6]
    \definition[位,名,个,些]{s.}{CEO; gerente geral; o mais alto executivo de uma empresa ou organização similar, que geralmente tem o poder de decidir políticas administrativas e de gestão}
  \end{Phonetics}
\end{Entry}

\begin{Entry}{总是}{9,9}{⼼,⽇}
  \begin{Phonetics}{总是}{zong3shi4}[][HSK 3]
    \definition{adv.}{sempre; indica como tem sido durante um determinado período de tempo; um determinado estado permanece inalterado | afinal; significa que, independentemente do que acontecer, haverá ou será um resultado}
  \end{Phonetics}
\end{Entry}

\begin{Entry}{总结}{9,9}{⼼,⽷}
  \begin{Phonetics}{总结}{zong3jie2}[][HSK 3]
    \definition[个,篇]{s.}{resumo; síntese; conclusão resumida}
    \definition{v.}{resumir; sumariar; sintetizar; analisar e estudar as experiências para chegar a conclusões}
  \end{Phonetics}
\end{Entry}

\begin{Entry}{总统}{9,9}{⼼,⽷}
  \begin{Phonetics}{总统}{zong3tong3}[][HSK 4]
    \definition*[个,位,名,家]{s.}{Presidente (de um país); Título dos líderes de determinadas repúblicas}
  \end{Phonetics}
\end{Entry}

\begin{Entry}{总值}{9,10}{⼼,⼈}
  \begin{Phonetics}{总值}{zong3zhi2}
    \definition{s.}{valor total; valor bruto; total geral; o valor monetário de algo que constitui o número ou a quantidade total é geralmente expresso como o preço de mercado calculado em termos da moeda de troca}
  \end{Phonetics}
\end{Entry}

\begin{Entry}{总监}{9,10}{⼼,⽫}
  \begin{Phonetics}{总监}{zong3jian1}[][HSK 6]
    \definition[名,位]{s.}{inspetor geral; inspetor"-chefe}
  \end{Phonetics}
\end{Entry}

\begin{Entry}{总站}{9,10}{⼼,⽴}
  \begin{Phonetics}{总站}{zong3zhan4}
    \definition{s.}{terminal}
  \end{Phonetics}
\end{Entry}

\begin{Entry}{总部}{9,10}{⼼,⾢}
  \begin{Phonetics}{总部}{zong3bu4}[][HSK 6]
    \definition{s.}{sede geral; escritório central}
  \end{Phonetics}
\end{Entry}

\begin{Entry}{总得}{9,11}{⼼,⼻}
  \begin{Phonetics}{总得}{zong3dei3}
    \definition{adv.}{prestes a; deve; tem que; está obrigado a}
    \definition{v.}{dever | precisar}
  \end{Phonetics}
\end{Entry}

\begin{Entry}{总理}{9,11}{⼼,⽟}
  \begin{Phonetics}{总理}{zong3li3}[][HSK 4]
    \definition*[个,位,名]{s.}{Primeiro"-Ministro do Conselho de Estado; Título do líder do Conselho de Estado da China | Título do chefe de governo em determinados países | Primeiro"-Ministro; Título de líderes de determinados partidos políticos | Título dos chefes de determinadas instituições e empresas nos velhos tempos}
    \definition{v.}{assumir a responsabilidade total}
  \end{Phonetics}
\end{Entry}

\begin{Entry}{总裁}{9,12}{⼼,⾐}
  \begin{Phonetics}{总裁}{zong3cai2}[][HSK 5]
    \definition[位,名,个]{s.}{presidente (de uma empresa); nomes de certos líderes de partidos políticos ou grandes empresas}
  \end{Phonetics}
\end{Entry}

\begin{Entry}{总量}{9,12}{⼼,⾥}
  \begin{Phonetics}{总量}{zong3liang4}[][HSK 6]
    \definition{s.}{capacidade total; quantidade bruta | valor total | total}
  \end{Phonetics}
\end{Entry}

\begin{Entry}{总数}{9,13}{⼼,⽁}
  \begin{Phonetics}{总数}{zong3shu4}[][HSK 5]
    \definition{s.}{soma; total; totalidade; inventário; número total; soma total}
  \end{Phonetics}
\end{Entry}

\begin{Entry}{总督}{9,13}{⼼,⽬}
  \begin{Phonetics}{总督}{zong3du1}
    \definition*[位,个]{s.}{Governador"-Geral | Governador | Vice"-Rei}
  \end{Phonetics}
\end{Entry}

\begin{Entry}{总算}{9,14}{⼼,⽵}
  \begin{Phonetics}{总算}{zong3suan4}[][HSK 5]
    \definition{adv.}{finalmente; por fim; indica que, após um longo período de tempo, um desejo finalmente se tornou realidade | suficiente; considerando tudo; no geral; considerando todos os aspectos; significa que, em geral, está tudo bem}
  \end{Phonetics}
\end{Entry}

\begin{Entry}{总额}{9,15}{⼼,⾴}
  \begin{Phonetics}{总额}{zong3'e2}[][HSK 7-9]
    \definition{s.}{valor total}
  \synonymref{总数}{zong3shu4}
  \end{Phonetics}
\end{Entry}

%%%%%%%%%% 恍 %%%%%%%%%%
\subsection*{恍}\addcontentsline{loh}{figure}{恍}

\begin{Entry}{恍}{9}{⼼}
  \begin{Phonetics}{恍}{huang3}
    \definition{adv.}{(junto com 如, 若, etc.) parecer; como se | de repente}
  \seealsoref{如}{ru2}
  \seealsoref{若}{ruo4}
  \end{Phonetics}
\end{Entry}

\begin{Entry}{恍然大悟}{9,12,3,10}{⼼,⽕,⼤,⼼}
  \begin{Phonetics}{恍然大悟}{huang3ran2-da4wu4}[][HSK 7-9]
    \definition{expr.}{de repente ver a luz; de repente perceber o que aconteceu; perceber de repente}
  \end{Phonetics}
\end{Entry}

%%%%%%%%%% 恒 %%%%%%%%%%
\subsection*{恒}\addcontentsline{loh}{figure}{恒}

\begin{Entry}{恒}{9}{⼼}
  \begin{Phonetics}{恒}{heng2}
    \definition*{s.}{Sobrenome: Heng}
    \definition{adj.}{permanente; duradouro | usual; comum; constante | usual; frequente; constante}
    \definition{s.}{perseverança; constância}
  \end{Phonetics}
\end{Entry}

\begin{Entry}{恒山}{9,3}{⼼,⼭}
  \begin{Phonetics}{恒山}{heng2shan1}
    \definition*{s.}{Monte Heng em Shanxi, montanha norte das Cinco Montanhas Sagradas (五岳) | Distrito de Hengshan na cidade de Jixi (鸡西), Heilongjiang}
  \seealsoref{鸡西}{ji1xi1}
  \seealsoref{五岳}{wu3yue4}
  \end{Phonetics}
\end{Entry}

\begin{Entry}{恒星系}{9,9,7}{⼼,⽇,⽷}
  \begin{Phonetics}{恒星系}{heng2xing1xi4}
    \definition{s.}{sistema estelar; galáxia}
  \end{Phonetics}
\end{Entry}

%%%%%%%%%% 恢 %%%%%%%%%%
\subsection*{恢}\addcontentsline{loh}{figure}{恢}

\begin{Entry}{恢}{9}{⼼}
  \begin{Phonetics}{恢}{hui1}
    \definition{adj.}{extenso; vasto | grande; ótimo}
    \definition{v.}{recuperar; restaurar; restabelecer}
  \end{Phonetics}
\end{Entry}

\begin{Entry}{恢复}{9,9}{⼼,⼢}
  \begin{Phonetics}{恢复}{hui1fu4}[][HSK 5]
    \definition{v.}{retomar; renovar; restaurar; voltar a | reviver; recuperar; reencontrar | restaurar; restabelecer; reabilitar; regenerar; ressurgir; restabelecer alguém em; recuperar o que foi perdido}
  \end{Phonetics}
\end{Entry}

%%%%%%%%%% 恤 %%%%%%%%%%
\subsection*{恤}\addcontentsline{loh}{figure}{恤}

\begin{Entry}{恤}{9}{⼼}
  \begin{Phonetics}{恤}{xu4}
    \definition{v.}{ter pena; simpatizar | dar alívio; compensar}
  \end{Phonetics}
\end{Entry}

%%%%%%%%%% 恨 %%%%%%%%%%
\subsection*{恨}\addcontentsline{loh}{figure}{恨}

\begin{Entry}{恨}{9}{⼼}
  \begin{Phonetics}{恨}{hen4}[][HSK 5]
    \definition{s.}{ódio; resentimento}
    \definition{v.}{odiar; ressentir"-se}
  \end{Phonetics}
\end{Entry}

\begin{Entry}{恨不得}{9,4,11}{⼼,⼀,⼻}
  \begin{Phonetics}{恨不得}{hen4bu5de5}[][HSK 7-9]
    \definition{v.}{estar muito ansioso para; querer poder (fazer algo); mal poder esperar para; expressa um desejo ansioso de realizar algo, geralmente usado para coisas que não podem realmente ser feitas}
  \end{Phonetics}
\end{Entry}

%%%%%%%%%% 恰 %%%%%%%%%%
\subsection*{恰}\addcontentsline{loh}{figure}{恰}

\begin{Entry}{恰}{9}{⼼}
  \begin{Phonetics}{恰}{qia4}
    \definition{adj.}{adequado; apropriado; oportuno}
    \definition{adv.}{na hora certa; com sorte; na medida certa; pontual; exatamente}
  \end{Phonetics}
\end{Entry}

\begin{Entry}{恰巧}{9,5}{⼼,⼯}
  \begin{Phonetics}{恰巧}{qia4qiao3}[][HSK 7-9]
    \definition{adv.}{por acaso; felizmente ou infelizmente; perfeitamente; por coincidência}
  \end{Phonetics}
\end{Entry}

\begin{Entry}{恰好}{9,6}{⼼,⼥}
  \begin{Phonetics}{恰好}{qia4hao3}[][HSK 6]
    \definition{adv.}{na medida certa; como a sorte quis}
  \end{Phonetics}
\end{Entry}

\begin{Entry}{恰如其分}{9,6,8,4}{⼼,⼥,⼋,⼑}
  \begin{Phonetics}{恰如其分}{qia4ru2-qi2fen4}[][HSK 7-9]
    \definition{expr.}{``Na medida certa.''; apropriado; agir ou falar de maneira diplomática}
  \end{Phonetics}
\end{Entry}

\begin{Entry}{恰当}{9,6}{⼼,⼹}
  \begin{Phonetics}{恰当}{qia4dang4}[][HSK 6]
    \definition{adj.}{adequado; apropriado; conveniente; apropriado; a linguagem ou abordagem é muito apropriada}
  \end{Phonetics}
\end{Entry}

\begin{Entry}{恰到好处}{9,8,6,5}{⼼,⼑,⼥,⼡}
  \begin{Phonetics}{恰到好处}{qia4dao4hao3chu4}[][HSK 7-9]
    \definition{expr.}{``Na medida certa.''; perfeito (para o propósito ou ocasião); significa que as palavras e ações de alguém atingiram o ponto mais apropriado}
  \end{Phonetics}
\end{Entry}

\begin{Entry}{恰恰}{9,9}{⼼,⼼}
  \begin{Phonetics}{恰恰}{qia4qia4}[][HSK 6]
    \definition{adv.}{justamente; exatamente; precisamente; bem na hora}
  \end{Phonetics}
\end{Entry}

\begin{Entry}{恰恰相反}{9,9,9,4}{⼼,⼼,⽬,⼜}
  \begin{Phonetics}{恰恰相反}{qia4qia4 xiang1fan3}[][HSK 7-9]
    \definition{expr.}{pelo contrário}
  \end{Phonetics}
\end{Entry}

%%%%%%%%%% 恼 %%%%%%%%%%
\subsection*{恼}\addcontentsline{loh}{figure}{恼}

\begin{Entry}{恼}{9}{⼼}
  \begin{Phonetics}{恼}{nao3}
    \definition{adj.}{infeliz; preocupado; aflito; angustiado}
    \definition{v.}{perturbar; irritar; incomodar | ficar com raiva; ficar irritado}
  \end{Phonetics}
\end{Entry}

\begin{Entry}{恼羞成怒}{9,10,6,9}{⼼,⽺,⼽,⼼}
  \begin{Phonetics}{恼羞成怒}{nao3xiu1-cheng2nu4}[][HSK 7-9]
    \definition{expr.}{ficar furioso de vergonha; ser envergonhado a ponto de ficar com raiva; sentir vergonha a ponto de ficar com raiva; irritação; ficar furioso por causa da humilhação}
  \end{Phonetics}
\end{Entry}

%%%%%%%%%% 恋 %%%%%%%%%%
\subsection*{恋}\addcontentsline{loh}{figure}{恋}

\begin{Entry}{恋}{10}{⼼}
  \begin{Phonetics}{恋}{lian4}
    \definition*{s.}{Sobrenome: Lian}
    \definition{v.}{amor (romântico) | ansiar por; sentir"-se apegado a | amar; apaixonar"-se por | não querendo se separar de; sentir sua falta para sempre; não suportar ficar separado}
  \end{Phonetics}
\end{Entry}

\begin{Entry}{恋恋不舍}{10,10,4,8}{⼼,⼼,⼀,⾆}
  \begin{Phonetics}{恋恋不舍}{lian4lian4-bu4she4}[][HSK 7-9]
    \definition{expr.}{afastar"-se de algo; relutar em se separar; sentir"-se apegado a algo; descrevendo a relutância em partir}
  \end{Phonetics}
\end{Entry}

\begin{Entry}{恋爱}{10,10}{⼼,⽖}
  \begin{Phonetics}{恋爱}{lian4'ai4}[][HSK 5]
    \definition[个,场,段]{s.}{namoro; afeto; amor romântico; ações que demonstram o amor mútuo}
    \definition{v.}{amar; estar apaixonado}
  \end{Phonetics}
\end{Entry}

%%%%%%%%%% 恐 %%%%%%%%%%
\subsection*{恐}\addcontentsline{loh}{figure}{恐}

\begin{Entry}{恐}{10}{⼼}
  \begin{Phonetics}{恐}{kong3}
    \definition{adv.}{talvez; provavelmente}
    \definition{v.}{temer; recear; ter medo de | ameaçar; aterrorizar; intimidar}
  \end{Phonetics}
\end{Entry}

\begin{Entry}{恐龙}{10,5}{⼼,⿓}
  \begin{Phonetics}{恐龙}{kong3long2}[][HSK 7-9]
    \definition[只,头]{s.}{dinossauro | garota feia (gíria da \emph{Internet}, ofensiva)}
  \end{Phonetics}
\end{Entry}

\begin{Entry}{恐吓}{10,6}{⼼,⼝}
  \begin{Phonetics}{恐吓}{kong3he4}[][HSK 7-9]
    \definition{v.}{ameaçar; assustar; intimidar; ameaçar alguém com palavras ou meios ameaçadores}
  \end{Phonetics}
\end{Entry}

\begin{Entry}{恐怕}{10,8}{⼼,⼼}
  \begin{Phonetics}{恐怕}{kong3pa4}[][HSK 3]
    \definition{adv.}{talvez; provavelmente; pode ser; expressa suposição; estimativa. | por medo de; expressar estimativa e preocupação}
    \definition{v.}{ter medo de; temer; recear}
  \end{Phonetics}
\end{Entry}

\begin{Entry}{恐怖}{10,8}{⼼,⼼}
  \begin{Phonetics}{恐怖}{kong3bu4}[][HSK 7-9]
    \definition[部]{adj.}{terrível; aterrador; horripilante; medo causado por ameaças à vida ou por presenciar violência ou derramamento de sangue | assustador; aterrorizante | terroristas; o comportamento ou os métodos utilizados são extremamente cruéis e perversos, causando choque e medo}
  \end{Phonetics}
\end{Entry}

\begin{Entry}{恐怖主义}{10,8,5,3}{⼼,⼼,⼂,⼂}
  \begin{Phonetics}{恐怖主义}{kong3bu4zhu3yi4}
    \definition{adj.}{terrorista}
    \definition{s.}{terrorismo}
  \end{Phonetics}
\end{Entry}

\begin{Entry}{恐惧}{10,11}{⼼,⼼}
  \begin{Phonetics}{恐惧}{kong3ju4}[][HSK 7-9]
    \definition{adj.}{assustado; com medo; muito assustado}
  \end{Phonetics}
\end{Entry}

\begin{Entry}{恐慌}{10,12}{⼼,⼼}
  \begin{Phonetics}{恐慌}{kong3huang1}[][HSK 7-9]
    \definition{adj.}{pânico; em pânico; pânico devido ao medo}
    \definition{s.}{pânico; medo}
  \end{Phonetics}
\end{Entry}

%%%%%%%%%% 恩 %%%%%%%%%%
\subsection*{恩}\addcontentsline{loh}{figure}{恩}

\begin{Entry}{恩}{10}{⼼}
  \begin{Phonetics}{恩}{en1}
    \definition*{s.}{Sobrenome: En}
    \definition{s.}{bondade; favor; graça; gentileza}
  \end{Phonetics}
\end{Entry}

\begin{Entry}{恩人}{10,2}{⼼,⼈}
  \begin{Phonetics}{恩人}{en1ren2}[][HSK 6]
    \definition{s.}{benfeitor; uma pessoa que ajudou significativamente alguém}
  \end{Phonetics}
\end{Entry}

\begin{Entry}{恩怨}{10,9}{⼼,⼼}
  \begin{Phonetics}{恩怨}{en1yuan4}[][HSK 7-9]
    \definition{s.}{sentimento de gratidão ou ressentimento (inimizade) | ressentimento; queixa; velhas contas}
  \end{Phonetics}
\end{Entry}

\begin{Entry}{恩情}{10,11}{⼼,⼼}
  \begin{Phonetics}{恩情}{en1qing2}[][HSK 7-9]
    \definition{s.}{amor; bondade; afeição profunda}
  \end{Phonetics}
\end{Entry}

\begin{Entry}{恩惠}{10,12}{⼼,⼼}
  \begin{Phonetics}{恩惠}{en1hui4}[][HSK 7-9]
    \definition[份]{s.}{favor; generosidade | bondade; graça; benefícios concedidos ou recebidos}
  \end{Phonetics}
\end{Entry}

\begin{Entry}{恩赐}{10,12}{⼼,⾙}
  \begin{Phonetics}{恩赐}{en1ci4}[][HSK 7-9]
    \definition{s.}{favor; caridade; esmola}
    \definition{v.}{conceder (favores, caridade, etc.); recompensar}
  \end{Phonetics}
\end{Entry}

%%%%%%%%%% 恭 %%%%%%%%%%
\subsection*{恭}\addcontentsline{loh}{figure}{恭}

\begin{Entry}{恭}{10}{⼼}
  \begin{Phonetics}{恭}{gong1}
    \definition{adj.}{respeitoso; reverente | educado}
  \end{Phonetics}
\end{Entry}

\begin{Entry}{恭维}{10,11}{⼼,⽷}
  \begin{Phonetics}{恭维}{gong1wei2}[][HSK 7-9]
    \definition{v.}{bajular; elogiar}
  \end{Phonetics}
\end{Entry}

\begin{Entry}{恭喜}{10,12}{⼼,⼝}
  \begin{Phonetics}{恭喜}{gong1xi3}[][HSK 7-9]
    \definition{v.}{parabenizar; uma maneira educada de parabenizar alguém por seu feliz evento}
  \end{Phonetics}
\end{Entry}

%%%%%%%%%% 息 %%%%%%%%%%
\subsection*{息}\addcontentsline{loh}{figure}{息}

\begin{Entry}{息}{10}{⼼}
  \begin{Phonetics}{息}{xi1}
    \definition*{s.}{Sobrenome: Xi}
    \definition{s.}{respiração; o ar que entra e sai durante a respiração | notícias; mensagem | juros | interesse | os filhos de alguém}
    \definition{v.}{cessar; parar | descansar | crescer; multiplicar; procriar}
  \synonymref{休}{xiu1}
  \antonymref{作}{zuo4}
  \end{Phonetics}
\end{Entry}

\begin{Entry}{息息相关}{10,10,9,6}{⼼,⼼,⽬,⼋}
  \begin{Phonetics}{息息相关}{xi1xi1-xiang1guan1}[][HSK 7-9]
    \definition{expr.}{estreitamente relacionado; intimamente relacionado; inextricavelmente ligado; interconectado; estreitamente conectado; ``A respiração está interligada.'', é uma metáfora para um relacionamento íntimo; estreitamente ligado; inter"-relacionado; extremamente relevante}
  \seealsoref{息息相通}{xi1xi1-xiang1tong1}
  \synonymref{息息相通}{xi1xi1-xiang1tong1}
  \end{Phonetics}
\end{Entry}

\begin{Entry}{息息相通}{10,10,9,10}{⼼,⼼,⽬,⾡}
  \begin{Phonetics}{息息相通}{xi1xi1-xiang1tong1}
    \definition{expr.}{mutuamente conectados; estreitamente ligado; inter"-relacionado; intimamente relacionado}
  \end{Phonetics}
\end{Entry}

%%%%%%%%%% 恳 %%%%%%%%%%
\subsection*{恳}\addcontentsline{loh}{figure}{恳}

\begin{Entry}{恳}{10}{⼼}
  \begin{Phonetics}{恳}{ken3}
    \definition{adj.}{sério; sincero | cordial; honesto}
    \definition{v.}{pedir; suplicar; implorar; rogar}
  \end{Phonetics}
\end{Entry}

\begin{Entry}{恳求}{10,7}{⼼,⽔}
  \begin{Phonetics}{恳求}{ken3qiu2}[][HSK 7-9]
    \definition{v.}{implorar; suplicar; rogar; solicitar encarecidamente}
  \end{Phonetics}
\end{Entry}

%%%%%%%%%% 恶 %%%%%%%%%%
\subsection*{恶}\addcontentsline{loh}{figure}{恶}

\begin{Entry}{恶}{10}{⼼}
  \begin{Phonetics}{恶}{e3}
    \definition{part.}{elementos formadores de palavras}
  \end{Phonetics}
  \begin{Phonetics}{恶}{e4}[][HSK 7-9]
    \definition{adj.}{feroz | ruim; maligno; perverso | vicioso | feio | grosseiro}
    \definition{s.}{mal; vício; crime | maldade; comportamento muito ruim; coisas criminosas}
  \antonymref{善}{shan4}
  \end{Phonetics}
  \begin{Phonetics}{恶}{wu1}
    \definition{interj.}{``Droga!''; ``Ah não!''; expressa surpresa}
    \definition{pron.}{como?; por que?; refere"-se a um lugar ou coisa; expressa uma pergunta retórica; equivalente a 何 ou 怎么}
  \seealsoref{何}{he2}
  \seealsoref{怎么}{zen3me5}
  \end{Phonetics}
  \begin{Phonetics}{恶}{wu4}
    \definition{v.}{não gostar; odiar; detestar; repugnar}
  \end{Phonetics}
\end{Entry}

\begin{Entry}{恶化}{10,4}{⼼,⼔}
  \begin{Phonetics}{恶化}{e4hua4}[][HSK 7-9]
    \definition{v.}{piorar; deteriorar; exacerbar | piorar a situação}
  \end{Phonetics}
\end{Entry}

\begin{Entry}{恶心}{10,4}{⼼,⼼}
  \begin{Phonetics}{恶心}{e3xin5}[][HSK 4]
    \definition{adj.}{nauseante; repugnante}
    \definition{s.}{náusea; repugnância}
    \definition{v.}{repugnar; ser nauseante; sentir"-se mal | envergonhar (deliberadamente)}
  \end{Phonetics}
  \begin{Phonetics}{恶心}{e4xin1}
    \definition{s.}{mau hábito; hábito vicioso; vício}
  \end{Phonetics}
\end{Entry}

\begin{Entry}{恶劣}{10,6}{⼼,⼒}
  \begin{Phonetics}{恶劣}{e4lie4}[][HSK 7-9]
    \definition{adj.}{mau; odioso; abominável; repugnante; desprezível; muito mau; muito ruim}
  \end{Phonetics}
\end{Entry}

\begin{Entry}{恶性}{10,8}{⼼,⼼}
  \begin{Phonetics}{恶性}{e4xing4}[][HSK 7-9]
    \definition{adj.}{maligno; pernicioso; vicioso | produzindo o mal | rápido (declínio) | descontrolada (inflação) | vicioso (círculo) | perverso}
  \antonymref{良性}{liang2xing4}
  \end{Phonetics}
\end{Entry}

\begin{Entry}{恶意}{10,13}{⼼,⼼}
  \begin{Phonetics}{恶意}{e4yi4}[][HSK 7-9]
    \definition[丝]{s.}{malícia; má vontade; más intenções}
  \end{Phonetics}
\end{Entry}

\begin{Entry}{恶霸}{10,21}{⼼,⾬}
  \begin{Phonetics}{恶霸}{e4ba4}
    \definition{s.}{tirano local; déspota local; valentão; pessoas más que se apoiam em forças reacionárias para dominar uma região e oprimir o povo}
  \synonymref{霸王}{ba4wang2}
  \antonymref{绅士}{shen1shi4}
  \end{Phonetics}
\end{Entry}

%%%%%%%%%% 悄 %%%%%%%%%%
\subsection*{悄}\addcontentsline{loh}{figure}{悄}

\begin{Entry}{悄}{10}{⼼}
  \begin{Phonetics}{悄}{qiao1}
    \definition{adj.}{quieto; silencioso}
  \end{Phonetics}
  \begin{Phonetics}{悄}{qiao3}
    \definition{adj.}{quieto; silencioso | triste; preocupado; aflito}
  \end{Phonetics}
\end{Entry}

\begin{Entry}{悄悄}{10,10}{⼼,⼼}
  \begin{Phonetics}{悄悄}{qiao1qiao1}[][HSK 5]
    \definition{adv.}{silenciosamente; em silêncio; aos sussuros; sem som ou em voz baixa; com o mínimo de ruído possível}
  \end{Phonetics}
\end{Entry}

%%%%%%%%%% 悔 %%%%%%%%%%
\subsection*{悔}\addcontentsline{loh}{figure}{悔}

\begin{Entry}{悔}{10}{⼼}
  \begin{Phonetics}{悔}{hui3}
    \definition{v.}{lamentar; arrepender"-se}
  \end{Phonetics}
\end{Entry}

\begin{Entry}{悔恨}{10,9}{⼼,⼼}
  \begin{Phonetics}{悔恨}{hui3hen4}[][HSK 7-9]
    \definition{v.}{arrepender"-se profundamente; estar amargamente arrependido}
  \end{Phonetics}
\end{Entry}

%%%%%%%%%% 悦 %%%%%%%%%%
\subsection*{悦}\addcontentsline{loh}{figure}{悦}

\begin{Entry}{悦}{10}{⼼}
  \begin{Phonetics}{悦}{yue4}
    \definition*{s.}{Sobrenome: Yue}
    \definition{adj.}{Linterário: feliz; satisfeito; encantado}
    \definition{v.}{agradar; deliciar}
  \end{Phonetics}
\end{Entry}

\begin{Entry}{悦耳}{10,6}{⼼,⽿}
  \begin{Phonetics}{悦耳}{yue4'er3}[][HSK 7-9]
    \definition{adj.}{agradável ao ouvido; de som doce; melodioso}
  \synonymref{动听}{dong4ting1}
  \synonymref{好听}{hao3ting1}
  \synonymref{顺耳}{shun4'er3}
  \antonymref{刺耳}{ci4'er3}
  \antonymref{难听}{nan2ting1}
  \antonymref{跑调}{pao3diao4}
  \end{Phonetics}
\end{Entry}

%%%%%%%%%% 悉 %%%%%%%%%%
\subsection*{悉}\addcontentsline{loh}{figure}{悉}

\begin{Entry}{悉}{11}{⼼}
  \begin{Phonetics}{悉}{xi1}
    \definition*{s.}{Sobrenome: Xi}
    \definition{adj.}{tudo; inteiro; total | detalhado}
    \definition{v.}{saber; aprender; ser informado de}
  \end{Phonetics}
\end{Entry}

\begin{Entry}{悉心}{11,4}{⼼,⼼}
  \begin{Phonetics}{悉心}{xi1xin1}
    \definition{v.}{dedicar toda a atenção; ter o máximo cuidado; usar toda a  força}
  \synonymref{精心}{jing1xin1}
  \synonymref{用心}{yong4/xin1}
  \antonymref{敷衍}{fu1yan3}
  \antonymref{马虎}{ma3hu5}
  \antonymref{疏忽}{shu1hu5}
  \antonymref{应付}{ying4fu5}
  \end{Phonetics}
\end{Entry}

\begin{Entry}{悉尼}{11,5}{⼼,⼫}
  \begin{Phonetics}{悉尼}{xi1ni2}
    \definition*{s.}{Sidney, capital de Nova Gales do Sul, Austrália}
  \end{Phonetics}
\end{Entry}

\begin{Entry}{悉数}{11,13}{⼼,⽁}
  \begin{Phonetics}{悉数}{xi1shu3}
    \definition{adv.}{Literário: listar; enumerar em detalhes; contar por completo}
  \end{Phonetics}
  \begin{Phonetics}{悉数}{xi1shu4}
    \definition{adv.}{todos | cada um | toda a soma}
  \synonymref{全部}{quan2bu4}
  \synonymref{全面}{quan2mian4}
  \synonymref{全体}{quan2ti3}
  \synonymref{所有}{suo3you3}
  \synonymref{统统}{tong3tong3}
  \synonymref{一切}{yi2qie4}
  \synonymref{整个}{zheng3ge4}
  \synonymref{总共}{zong3gong4}
  \end{Phonetics}
\end{Entry}

%%%%%%%%%% 悠 %%%%%%%%%%
\subsection*{悠}\addcontentsline{loh}{figure}{悠}

\begin{Entry}{悠}{11}{⼼}
  \begin{Phonetics}{悠}{you1}
    \definition{adj.}{velho; demorado; remoto (no tempo ou no espaço) | tranquilo; relaxado; pensativo}
    \definition{v.}{balançar}
  \end{Phonetics}
\end{Entry}

\begin{Entry}{悠久}{11,3}{⼼,⼃}
  \begin{Phonetics}{悠久}{you1jiu3}[][HSK 7-9]
    \definition{adj.}{longo; ancestral; de longa data; descreve uma história, cultura ou outro período histórico que tenha durado muito tempo}
  \synonymref{长久}{chang2jiu3}
  \synonymref{长远}{chang2yuan3}
  \synonymref{持久}{chi2jiu3}
  \synonymref{好久}{hao3jiu3}
  \synonymref{漫长}{man4chang2}
  \synonymref{深远}{shen1yuan3}
  \synonymref{永久}{yong3jiu3}
  \synonymref{永远}{yong3yuan3}
  \antonymref{短促}{duan3cu4}
  \antonymref{短暂}{duan3zan4}
  \end{Phonetics}
\end{Entry}

\begin{Entry}{悠闲}{11,7}{⼼,⾨}
  \begin{Phonetics}{悠闲}{you1xian2}[][HSK 7-9]
    \definition{adj.}{descontraído e despreocupado; um estado de lazer e contentamento}
  \antonymref{匆忙}{cong1mang2}
  \antonymref{繁忙}{fan2mang2}
  \antonymref{忙碌}{mang2lu4}
  \end{Phonetics}
\end{Entry}

%%%%%%%%%% 患 %%%%%%%%%%
\subsection*{患}\addcontentsline{loh}{figure}{患}

\begin{Entry}{患}{11}{⼼}
  \begin{Phonetics}{患}{huan4}[][HSK 7-9]
    \definition*{s.}{Sobrenome: Huan}
    \definition{s.}{perigo; problema; desastre; flagelo | preocupação; ansiedade}
    \definition{v.}{contrair (doença); sofrer de}
  \end{Phonetics}
\end{Entry}

\begin{Entry}{患有}{11,6}{⼼,⽉}
  \begin{Phonetics}{患有}{huan4you3}[][HSK 7-9]
    \definition{v.}{sofrer de; refere"-se a alguém que sofre de uma doença ou condição específica}
  \end{Phonetics}
\end{Entry}

\begin{Entry}{患者}{11,8}{⼼,⽼}
  \begin{Phonetics}{患者}{huan4zhe3}[][HSK 6]
    \definition[个,位,名]{s.}{paciente; sofredor; pessoas com certas doenças}
  \end{Phonetics}
\end{Entry}

\begin{Entry}{患病}{11,10}{⼼,⽧}
  \begin{Phonetics}{患病}{huan4bing4}[][HSK 7-9]
    \definition{v.}{estar doente; ficar doente; adoecer; sofrer de uma doença}
  \end{Phonetics}
\end{Entry}

%%%%%%%%%% 您 %%%%%%%%%%
\subsection*{您}\addcontentsline{loh}{figure}{您}

\begin{Entry}{您}{11}{⼼}
  \begin{Phonetics}{您}{nin2}[][HSK 1]
    \definition{pron.}{você; a forma de tratamento respeitosa da segunda pessoa do singular 你}
  \seealsoref{你}{ni3}
  \synonymref{你}{ni3}
  \end{Phonetics}
\end{Entry}

%%%%%%%%%% 悬 %%%%%%%%%%
\subsection*{悬}\addcontentsline{loh}{figure}{悬}

\begin{Entry}{悬}{11}{⼼}
  \begin{Phonetics}{悬}{xuan2}[][HSK 6]
    \definition{adj.}{pendente; não resolvido; sem nenhum resultado | distante; a distância é grande; a diferença é grande | Dialeto: perigoso}
    \definition{v.}{pendurar; suspender | levantar; elevar | sentir"-se ansioso; ser solícito | imaginar}
  \end{Phonetics}
\end{Entry}

\begin{Entry}{悬心吊胆}{11,4,6,9}{⼼,⼼,⼝,⾁}
  \begin{Phonetics}{悬心吊胆}{xuan2xin1-diao4dan3}
    \definition{expr.}{de tirar o fôlego; estar ansioso e com medo; estar à beira de um ataque de nervos; ter o coração na boca; estar em suspense}
  \synonymref{提心吊胆}{ti2xin1-diao4dan3}
  \end{Phonetics}
\end{Entry}

\begin{Entry}{悬念}{11,8}{⼼,⼼}
  \begin{Phonetics}{悬念}{xuan2nian4}[][HSK 7-9]
    \definition[处]{s.}{suspense (sentido à medida que uma história, peça teatral, etc., se desenvolve até um clímax); o sentimento de preocupação com o desenvolvimento da história e o destino dos personagens ao apreciar uma peça de teatro, um filme ou outra obra de arte}
    \definition{v.}{preocupar"-se com; ficar pensando em; manter em suspense}
  \synonymref{担心}{dan1/xin1}
  \synonymref{惦记}{dian4ji4}
  \synonymref{顾虑}{gu4lv4}
  \synonymref{挂念}{gua4nian4}
  \synonymref{怀念}{huai2nian4}
  \synonymref{牵挂}{qian1gua4}
  \synonymref{思念}{si1nian4}
  \synonymref{想念}{xiang3nian4}
  \antonymref{放心}{fang4/xin1}
  \end{Phonetics}
\end{Entry}

\begin{Entry}{悬挂}{11,9}{⼼,⼿}
  \begin{Phonetics}{悬挂}{xuan2gua4}[][HSK 7-9]
    \definition{v.}{pendurar; pender; suspender; prender um objeto em um ou mais pontos em algum lugar com a ajuda de uma corda, gancho, prego, etc.}
  \end{Phonetics}
\end{Entry}

\begin{Entry}{悬殊}{11,10}{⼼,⽍}
  \begin{Phonetics}{悬殊}{xuan2shu1}[][HSK 7-9]
    \definition{adj.}{vastamente diferentes; polos opostos}
  \synonymref{差距}{cha1ju4}
  \antonymref{相当}{xiang1dang1}
  \end{Phonetics}
\end{Entry}

\begin{Entry}{悬崖}{11,11}{⼼,⼭}
  \begin{Phonetics}{悬崖}{xuan2ya2}[][HSK 7-9]
    \definition[处]{s.}{penhasco; precipício; escarpa}
  \end{Phonetics}
\end{Entry}

%%%%%%%%%% 悼 %%%%%%%%%%
\subsection*{悼}\addcontentsline{loh}{figure}{悼}

\begin{Entry}{悼}{11}{⼼}
  \begin{Phonetics}{悼}{dao4}
    \definition*{s.}{Sobrenome: Dao}
    \definition{v.}{lamentar; expressar pesar}
  \end{Phonetics}
\end{Entry}

\begin{Entry}{悼念}{11,8}{⼼,⼼}
  \begin{Phonetics}{悼念}{dao4nian4}[][HSK 7-9]
    \definition{v.}{lamentar; lamentar"-se por | lamentar por; expressar pesar}
  \end{Phonetics}
\end{Entry}

%%%%%%%%%% 情 %%%%%%%%%%
\subsection*{情}\addcontentsline{loh}{figure}{情}

\begin{Entry}{情}{11}{⼼}
  \begin{Phonetics}{情}{qing2}[][HSK 7-9]
    \definition{s.}{sentimento; afeição | amor; paixão | paixão sexual; luxúria | favor; gentileza | situação; circunstâncias; condição | razão; sentido | sensibilidades; sentimentos}
  \end{Phonetics}
\end{Entry}

\begin{Entry}{情人}{11,2}{⼼,⼈}
  \begin{Phonetics}{情人}{qing2ren2}[][HSK 7-9]
    \definition[对,个,位]{s.}{amante; namorado(a) | concubina; amante}
  \end{Phonetics}
\end{Entry}

\begin{Entry}{情不自禁}{11,4,6,13}{⼼,⼀,⾃,⽰}
  \begin{Phonetics}{情不自禁}{qing2bu2zi4jin1}[][HSK 7-9]
    \definition{expr.}{``Não consigo ajudar.''; não conseguir se conter; não conseguir evitar (fazer algo); ser tomado por um impulso repentino de; dominado pela emoção; enfatizando o controle total sobre as próprias emoções}
  \end{Phonetics}
\end{Entry}

\begin{Entry}{情节}{11,5}{⼼,⾋}
  \begin{Phonetics}{情节}{qing2jie2}[][HSK 5]
    \definition[个,段]{s.}{enredo; trama; desenrolar específico dos acontecimentos | circunstância; detalhes do crime ou erro | enredo; roteiro; refere"-se especificamente ao processo de desenvolvimento e evolução dos conflitos e contradições em obras literárias narrativas}
  \end{Phonetics}
\end{Entry}

\begin{Entry}{情况}{11,7}{⼼,⼎}
  \begin{Phonetics}{情况}{qing2kuang4}[][HSK 3]
    \definition[种,个,些]{s.}{condição; situação; circunstâncias; estado das coisas | mudanças notáveis e impactantes}
  \end{Phonetics}
\end{Entry}

\begin{Entry}{情形}{11,7}{⼼,⼺}
  \begin{Phonetics}{情形}{qing2xing5}[][HSK 5]
    \definition[个,种]{s.}{situação; condição; circunstâncias; estado de coisas; a situação específica das coisas}
  \end{Phonetics}
\end{Entry}

\begin{Entry}{情怀}{11,7}{⼼,⼼}
  \begin{Phonetics}{情怀}{qing2huai2}[][HSK 7-9]
    \definition{s.}{sentimentos; um estado de espírito que contém uma determinada emoção}
  \end{Phonetics}
\end{Entry}

\begin{Entry}{情报}{11,7}{⼼,⼿}
  \begin{Phonetics}{情报}{qing2bao4}[][HSK 7-9]
    \definition[个,份]{s.}{inteligência; informação; notícias e reportagens sobre determinada situação são frequentemente classificadas como confidenciais}
  \end{Phonetics}
\end{Entry}

\begin{Entry}{情侣}{11,8}{⼼,⼈}
  \begin{Phonetics}{情侣}{qing2lv3}[][HSK 7-9]
    \definition[对,双,群]{s.}{amantes; namorados; um casal apaixonado ou um deles}
  \end{Phonetics}
\end{Entry}

\begin{Entry}{情结}{11,9}{⼼,⽷}
  \begin{Phonetics}{情结}{qing2jie2}[][HSK 7-9]
    \definition{s.}{complexidade; a turbulência emocional em meu coração; um certo sentimento que frequentemente persiste em minha mente}
  \end{Phonetics}
\end{Entry}

\begin{Entry}{情调}{11,10}{⼼,⾔}
  \begin{Phonetics}{情调}{qing2diao4}[][HSK 7-9]
    \definition[出]{s.}{sentimento; apelo emocional; tom afetivo; o estilo expresso através de pensamentos e sentimentos; a natureza das coisas que podem evocar diversas emoções diferentes nas pessoas}
  \end{Phonetics}
\end{Entry}

\begin{Entry}{情谊}{11,10}{⼼,⾔}
  \begin{Phonetics}{情谊}{qing2yi4}[][HSK 7-9]
    \definition{s.}{amizade; sentimentos amigáveis; emoções amistosas; os sentimentos de carinho e amor entre pessoas}
  \end{Phonetics}
\end{Entry}

\begin{Entry}{情绪}{11,11}{⼼,⽷}
  \begin{Phonetics}{情绪}{qing2xu4}[][HSK 6]
    \definition[种,片,股,丝]{s.}{mau humor; depressão; um sentimento ruim no coração, especialmente um estado mental desagradável quando se sente injusto | emoção; humor; moral; sentimento; o estado mental de uma pessoa ao longo de um período de tempo}
  \end{Phonetics}
\end{Entry}

\begin{Entry}{情景}{11,12}{⼼,⽇}
  \begin{Phonetics}{情景}{qing2jing3}[][HSK 4]
    \definition[个,幕,种]{s.}{cena; vista; circunstâncias}
  \end{Phonetics}
\end{Entry}

\begin{Entry}{情感}{11,13}{⼼,⼼}
  \begin{Phonetics}{情感}{qing2gan3}[][HSK 3]
    \definition[份]{s.}{emoção; sentimento | afeição; apego; reações psicológicas positivas ou negativas a estímulos externos, como gosto, raiva, tristeza, medo, amor, nojo, etc.}
  \end{Phonetics}
\end{Entry}

\begin{Entry}{情愿}{11,14}{⼼,⽕}
  \begin{Phonetics}{情愿}{qing2yuan4}[][HSK 7-9]
    \definition{v.}{estar disposto a; estar genuinamente disposto a fazer algo que outros não estão dispostos a fazer}
  \end{Phonetics}
\end{Entry}

%%%%%%%%%% 惊 %%%%%%%%%%
\subsection*{惊}\addcontentsline{loh}{figure}{惊}

\begin{Entry}{惊}{11}{⼼}
  \begin{Phonetics}{惊}{jing1}[][HSK 7-9]
    \definition{v.}{assustar; ficar assustado; ficar nervoso devido a estímulo repentino; ficar com medo | surpreender; chocar; alarmar}
  \end{Phonetics}
\end{Entry}

\begin{Entry}{惊人}{11,2}{⼼,⼈}
  \begin{Phonetics}{惊人}{jing1ren2}[][HSK 6]
    \definition{adj.}{surpreso; espantado; atônito; surpreendente}
  \end{Phonetics}
\end{Entry}

\begin{Entry}{惊天动地}{11,4,6,6}{⼼,⼤,⼒,⼟}
  \begin{Phonetics}{惊天动地}{jing1tian1-dong4di4}[][HSK 7-9]
    \definition{expr.}{que abala os céus e a terra; que faz a terra tremer; assustar os céus e mover a terra; sacudir os céus e assustar a terra; que abala o mundo; de proporções sísmicas; de impacto mundial}
  \end{Phonetics}
\end{Entry}

\begin{Entry}{惊心动魄}{11,4,6,14}{⼼,⼼,⼒,⿁}
  \begin{Phonetics}{惊心动魄}{jing1xin1-dong4po4}[][HSK 7-9]
    \definition{expr.}{comovente; profundamente impactante; ficar apavorado (horror); de tirar o fôlego; de arrepiar os cabelos; emocionante; fazer o coração de alguém disparar; abalar alguém profundamente; deixar alguém sem fôlego}
  \end{Phonetics}
\end{Entry}

\begin{Entry}{惊叹}{11,5}{⼼,⼝}
  \begin{Phonetics}{惊叹}{jing1tan4}[][HSK 7-9]
    \definition{v.}{maravilhar"-se com; admirar"-se com; exclamar (com admiração)}
  \end{Phonetics}
\end{Entry}

\begin{Entry}{惊异}{11,6}{⼼,⼶}
  \begin{Phonetics}{惊异}{jing1yi4}
    \definition{adj.}{surpreso; admirado; estupefato; atônito}
  \end{Phonetics}
\end{Entry}

\begin{Entry}{惊讶}{11,6}{⼼,⾔}
  \begin{Phonetics}{惊讶}{jing1ya4}[][HSK 7-9]
    \definition{adj.}{surpreso; admirado; estupefato; atônito; sentindo"-me surpreso e estranho}
  \end{Phonetics}
\end{Entry}

\begin{Entry}{惊呆}{11,7}{⼼,⼝}
  \begin{Phonetics}{惊呆}{jing1dai1}
    \definition{adj.}{atordoado; estupefato; chocado}
  \end{Phonetics}
\end{Entry}

\begin{Entry}{惊奇}{11,8}{⼼,⼤}
  \begin{Phonetics}{惊奇}{jing1qi2}[][HSK 7-9]
    \definition{v.}{maravilhar"-se; ficar surpreso; ficar boquiaberto}
  \end{Phonetics}
\end{Entry}

\begin{Entry}{惊诧}{11,8}{⼼,⾔}
  \begin{Phonetics}{惊诧}{jing1cha4}[][HSK 7-9]
    \definition{adj.}{surpreso; admirado; estupefato | muito surpreso}
  \end{Phonetics}
\end{Entry}

\begin{Entry}{惊险}{11,9}{⼼,⾩}
  \begin{Phonetics}{惊险}{jing1xian3}[][HSK 7-9]
    \definition{adj.}{emocionante; de tirar o fôlego; alarmantemente perigoso}
  \end{Phonetics}
\end{Entry}

\begin{Entry}{惊喜}{11,12}{⼼,⼝}
  \begin{Phonetics}{惊喜}{jing1xi3}[][HSK 6]
    \definition{s.}{boa surpresa; agradavelmente surpreso}
  \end{Phonetics}
\end{Entry}

\begin{Entry}{惊慌}{11,12}{⼼,⼼}
  \begin{Phonetics}{惊慌}{jing1huang1}[][HSK 7-9]
    \definition{adj.}{assustado; alarmado; amedrontado; atemorizado; em pânico}
  \end{Phonetics}
\end{Entry}

\begin{Entry}{惊慌失措}{11,12,5,11}{⼼,⼼,⼤,⼿}
  \begin{Phonetics}{惊慌失措}{jing1huang1-shi1cuo4}[][HSK 7-9]
    \definition{expr.}{apavorado; tomado pelo pânico; em pânico; perder a cabeça de medo; no Capítulo 11 da Parte 4 de ``O Oriente'', de Wei Wei: ``Guo Xiang e seus homens lançaram um ataque feroz, e o inimigo, em pânico, esqueceu"-se de resistir, preocupando"-se apenas em subir nos tanques.''; confusão aterrorizante; ficar apavorado; ficar em pânico}
  \end{Phonetics}
\end{Entry}

\begin{Entry}{惊醒}{11,16}{⼼,⾣}
  \begin{Phonetics}{惊醒}{jing1xing3}[][HSK 7-9]
    \definition{v.}{acordar sobressaltado; despertado pelo susto; despertado pelo choque}
  \end{Phonetics}
\end{Entry}

%%%%%%%%%% 惋 %%%%%%%%%%
\subsection*{惋}\addcontentsline{loh}{figure}{惋}

\begin{Entry}{惋}{11}{⼼}
  \begin{Phonetics}{惋}{wan3}
    \definition{s.}{Literário: suspiro}
    \definition{v.}{Literário: suspirar}
  \end{Phonetics}
\end{Entry}

\begin{Entry}{惋惜}{11,11}{⼼,⼼}
  \begin{Phonetics}{惋惜}{wan3xi1}[][HSK 7-9]
    \definition{v.}{lamentar; sentir que é uma pena; sentir pena de alguém; estar arrependido; expressar simpatia e pena pelos infortúnios das pessoas ou por mudanças insatisfatórias nas coisas}
  \synonymref{可惜}{ke3xi1}
  \synonymref{怜惜}{lian2xi1}
  \synonymref{无奈}{wu2nai4}
  \synonymref{遗憾}{yi2han4}
  \antonymref{庆幸}{qing4xing4}
  \end{Phonetics}
\end{Entry}

%%%%%%%%%% 惦 %%%%%%%%%%
\subsection*{惦}\addcontentsline{loh}{figure}{惦}

\begin{Entry}{惦}{11}{⼼}
  \begin{Phonetics}{惦}{dian4}
    \definition{v.}{lembrar com preocupação; estar preocupado com; continuar pensando sobre; ficar pensando em alguém ou em alguma coisa e se preocupar com eles; sentir falta deles}
  \end{Phonetics}
\end{Entry}

\begin{Entry}{惦记}{11,5}{⼼,⾔}
  \begin{Phonetics}{惦记}{dian4ji4}[][HSK 7-9]
    \definition{s.}{lembrar com preocupação; estar preocupado com; continuar pensando sobre; (sobre uma pessoa ou coisa) continuar pensando nisso e não deixar passar}
  \end{Phonetics}
\end{Entry}

%%%%%%%%%% 惭 %%%%%%%%%%
\subsection*{惭}\addcontentsline{loh}{figure}{惭}

\begin{Entry}{惭}{11}{⼼}
  \begin{Phonetics}{惭}{can2}
    \definition{adj.}{envergonhado}
    \definition{s.}{vergonha}
    \definition{v.}{sentir vergonha}
  \end{Phonetics}
\end{Entry}

\begin{Entry}{惭愧}{11,12}{⼼,⼼}
  \begin{Phonetics}{惭愧}{can2kui4}[][HSK 7-9]
    \definition{adj.}{envergonhado; sentir"-se inseguro por ter deficiências, fazer algo errado ou não cumprir responsabilidades}
  \end{Phonetics}
\end{Entry}

%%%%%%%%%% 惯 %%%%%%%%%%
\subsection*{惯}\addcontentsline{loh}{figure}{惯}

\begin{Entry}{惯}{11}{⼼}
  \begin{Phonetics}{惯}{guan4}[][HSK 7-9]
    \definition{adj.}{habitual; costumeiro; usual | incorrigível; endurecido}
    \definition{v.}{estar acostumado a; ter o hábito de | mimar; estragar}
  \end{Phonetics}
\end{Entry}

\begin{Entry}{惯例}{11,8}{⼼,⼈}
  \begin{Phonetics}{惯例}{guan4li4}[][HSK 7-9]
    \definition[个]{s.}{rotina; convenção; prática usual; prática habitual | precedente; embora não haja nenhuma disposição explícita na lei, há práticas que foram implementadas no passado e podem ser imitadas}
  \end{Phonetics}
\end{Entry}

\begin{Entry}{惯性}{11,8}{⼼,⼼}
  \begin{Phonetics}{惯性}{guan4xing4}[][HSK 7-9]
    \definition{s.}{Física: inércia; a força da inércia}
  \end{Phonetics}
\end{Entry}

%%%%%%%%%% 惑 %%%%%%%%%%
\subsection*{惑}\addcontentsline{loh}{figure}{惑}

\begin{Entry}{惑}{12}{⼼}
  \begin{Phonetics}{惑}{huo4}
    \definition{v.}{ficar confuso; ficar perplexo | iludir; enganar; confundir}
  \end{Phonetics}
\end{Entry}

\begin{Entry}{惑星}{12,9}{⼼,⽇}
  \begin{Phonetics}{惑星}{huo4xing1}
    \definition{s.}{planeta; também escrito como 行星}
  \seealsoref{行星}{xing2xing1}
  \end{Phonetics}
\end{Entry}

%%%%%%%%%% 惩 %%%%%%%%%%
\subsection*{惩}\addcontentsline{loh}{figure}{惩}

\begin{Entry}{惩}{12}{⼼}
  \begin{Phonetics}{惩}{cheng2}
    \definition{v.}{receber ou dar aviso | punir; penalizar}
  \end{Phonetics}
\end{Entry}

\begin{Entry}{惩处}{12,5}{⼼,⼡}
  \begin{Phonetics}{惩处}{cheng2chu3}[][HSK 7-9]
    \definition{v.}{penalizar; punir | punir; administrar justiça}
  \end{Phonetics}
\end{Entry}

\begin{Entry}{惩罚}{12,9}{⼼,⽹}
  \begin{Phonetics}{惩罚}{cheng2fa2}[][HSK 7-9]
    \definition[次,种]{s.}{punição; o ato ou método de punição}
    \definition{v.}{punir (severamente); penalizar}
  \end{Phonetics}
\end{Entry}

%%%%%%%%%% 惴 %%%%%%%%%%
\subsection*{惴}\addcontentsline{loh}{figure}{惴}

\begin{Entry}{惴}{12}{⼼}
  \begin{Phonetics}{惴}{zhui4}
    \definition{adj.}{ansioso e medroso; Literário: é ao mesmo tempo preocupante e assustador}
  \end{Phonetics}
\end{Entry}

%%%%%%%%%% 惶 %%%%%%%%%%
\subsection*{惶}\addcontentsline{loh}{figure}{惶}

\begin{Entry}{惶}{12}{⼼}
  \begin{Phonetics}{惶}{huang2}
    \definition{adj.}{cheio de medo; assustado}
    \definition{s.}{medo; pânico}
    \definition{v.}{temer}
  \end{Phonetics}
\end{Entry}

\begin{Entry}{惶恐}{12,10}{⼼,⼼}
  \begin{Phonetics}{惶恐}{huang2kong3}
    \definition{adj.}{aterrorizado; em pânico; petrificado | inquieto; apreensivo}
  \end{Phonetics}
\end{Entry}

%%%%%%%%%% 惹 %%%%%%%%%%
\subsection*{惹}\addcontentsline{loh}{figure}{惹}

\begin{Entry}{惹}{12}{⼼}
  \begin{Phonetics}{惹}{re3}[][HSK 7-9]
    \definition{v.}{provocar; convidar ou pedir (algo indesejável); (características de uma pessoa ou coisa) evocar sentimentos de amor ou ódio | ofender; provocar; instigar; (palavras e ações) tocar na outra pessoa | atrair; causar (algo ruim)}
  \end{Phonetics}
\end{Entry}

%%%%%%%%%% 愉 %%%%%%%%%%
\subsection*{愉}\addcontentsline{loh}{figure}{愉}

\begin{Entry}{愉}{12}{⼼}
  \begin{Phonetics}{愉}{yu2}
    \definition{adj.}{satisfeito; feliz; alegre}
  \end{Phonetics}
\end{Entry}

\begin{Entry}{愉快}{12,7}{⼼,⼼}
  \begin{Phonetics}{愉快}{yu2kuai4}[][HSK 6]
    \definition{adj.}{feliz; alegre; de bom humor, muito feliz}
  \end{Phonetics}
\end{Entry}

%%%%%%%%%% 愣 %%%%%%%%%%
\subsection*{愣}\addcontentsline{loh}{figure}{愣}

\begin{Entry}{愣}{12}{⼼}
  \begin{Phonetics}{愣}{leng4}[][HSK 7-9]
    \definition{adj.}{precipitado; rude; imprudente; temerário; falar e agir sem considerar as consequências}
    \definition{adv.}{insistir em algo de forma veemente, sem provas ou fundamentos}
    \definition{v.}{ficar entorpecido; ficar em branco; ficar distraído; ficar atordoado; ficar desatento; ficar perdido em pensamentos}
  \end{Phonetics}
\end{Entry}

%%%%%%%%%% 愤 %%%%%%%%%%
\subsection*{愤}\addcontentsline{loh}{figure}{愤}

\begin{Entry}{愤}{12}{⼼}
  \begin{Phonetics}{愤}{fen4}
    \definition{s.}{raiva; indignação; ressentimento; exasperação}
    \definition{v.}{ressentir"-se; ficar indignado; ficar com raiva}
  \end{Phonetics}
\end{Entry}

\begin{Entry}{愤世嫉俗}{12,5,13,9}{⼼,⼀,⼥,⼈}
  \begin{Phonetics}{愤世嫉俗}{fen4shi4ji2su2}
    \definition{expr.}{detestar o mundo e seus costumes; ser cínico; ficar amargurado}
  \end{Phonetics}
\end{Entry}

\begin{Entry}{愤怒}{12,9}{⼼,⼼}
  \begin{Phonetics}{愤怒}{fen4nu4}[][HSK 6]
    \definition{adj.}{zangado; enraivecido; iracundo; furioso; emocionalmente agitado por extrema insatisfação}
  \end{Phonetics}
\end{Entry}

%%%%%%%%%% 慌 %%%%%%%%%%
\subsection*{慌}\addcontentsline{loh}{figure}{慌}

\begin{Entry}{慌}{12}{⼼}
  \begin{Phonetics}{慌}{huang1}[][HSK 5]
    \definition{adj.}{agitado; perturbado; confuso; que inspira terror}
    \definition{v.}{estar em estado de pânico; ficar com medo; ficar nervoso | estar com pressa}
  \end{Phonetics}
\end{Entry}

\begin{Entry}{慌忙}{12,6}{⼼,⼼}
  \begin{Phonetics}{慌忙}{huang1mang2}[][HSK 5]
    \definition{adj.}{apressado; afobado; com muita pressa}
    \definition{adv.}{apressadamente}
  \end{Phonetics}
\end{Entry}

\begin{Entry}{慌乱}{12,7}{⼼,⼄}
  \begin{Phonetics}{慌乱}{huang1luan4}[][HSK 7-9]
    \definition{adj.}{agitado; alarmado e confuso; em pânico e ocupado}
  \end{Phonetics}
\end{Entry}

\begin{Entry}{慌张}{12,7}{⼼,⼸}
  \begin{Phonetics}{慌张}{huang1zhang1}[][HSK 7-9]
    \definition{adj.}{em pânico; agitado; perturbado; confuso}
  \end{Phonetics}
\end{Entry}

%%%%%%%%%% 想 %%%%%%%%%%
\subsection*{想}\addcontentsline{loh}{figure}{想}

\begin{Entry}{想}{13}{⼼}
  \begin{Phonetics}{想}{xiang3}[][HSK 1]
    \definition{v.}{pensar; ponderar; refletir | supor; contar; considerar; pensar; estimar | querer; gostaria de; sentir vontade (de fazer algo) | lembrar com saudade; sentir falta}
  \seealsoref{想念}{xiang3nian4}
  \seealsoref{要}{yao4}
  \synonymref{记}{ji4}
  \synonymref{念}{nian4}
  \synonymref{思}{si1}
  \synonymref{希}{xi1}
  \synonymref{想念}{xiang3nian4}
  \synonymref{要}{yao4}
  \end{Phonetics}
\end{Entry}

\begin{Entry}{想不到}{13,4,8}{⼼,⼀,⼑}
  \begin{Phonetics}{想不到}{xiang3bu2dao4}[][HSK 6]
    \definition{adj.}{inesperado; imprevisto}
  \end{Phonetics}
\end{Entry}

\begin{Entry}{想方设法}{13,4,6,8}{⼼,⽅,⾔,⽔}
  \begin{Phonetics}{想方设法}{xiang3fang1-she4fa3}[][HSK 7-9]
    \definition{expr.}{fazer tudo o que for possível; tentar todos os meios; tentar por bem ou por mal}
  \synonymref{千方百计}{qian1fang1-bai3ji4}
  \end{Phonetics}
\end{Entry}

\begin{Entry}{想到}{13,8}{⼼,⼑}
  \begin{Phonetics}{想到}{xiang3dao4}[][HSK 2]
    \definition{v.}{pensar em; trazer à mente; ter no coração; ter uma ideia (na mente); ter uma ideia (no coração)}
  \seealsoref{料到}{liao4dao4}
  \synonymref{料到}{liao4dao4}
  \end{Phonetics}
\end{Entry}

\begin{Entry}{想念}{13,8}{⼼,⼼}
  \begin{Phonetics}{想念}{xiang3nian4}[][HSK 4]
    \definition{v.}{sentir falta; pensar em; lembrar com carinho; ficar doente por; desejar ver novamente; lembrar com saudade}
  \end{Phonetics}
\end{Entry}

\begin{Entry}{想法}{13,8}{⼼,⽔}
  \begin{Phonetics}{想法}{xiang3fa5}[][HSK 2]
    \definition[种]{s.}{ideia; opinião; pensamento; noção; o que alguém tem em mente; visões e opiniões sobre alguém ou algo obtidas através do pensamento}
    \definition{s.}{maneira de pensar | opinião | noção}
    \definition{v.}{tentar; pensar em uma maneira (de fazer algo); fazer o que puder; encontrar um jeito}
  \seealsoref{观念}{guan1nian4}
  \seealsoref{看法}{kan4fa5}
  \synonymref{点子}{dian3zi5}
  \synonymref{观点}{guan1dian3}
  \synonymref{观念}{guan1nian4}
  \synonymref{见解}{jian4jie3}
  \synonymref{念头}{nian4tou5}
  \synonymref{设法}{she4fa3}
  \synonymref{思想}{si1xiang3}
  \synonymref{主张}{zhu3zhang1}
  \synonymref{主意}{zhu3yi5}
  \end{Phonetics}
\end{Entry}

\begin{Entry}{想起}{13,10}{⼼,⾛}
  \begin{Phonetics}{想起}{xiang3qi3}[][HSK 2]
    \definition{v.}{recordar; lembrar; pensar em; trazer à mente; cruzar pelos pensamentos de alguém; passar pelo pensamento de alguém}
  \end{Phonetics}
\end{Entry}

\begin{Entry}{想象}{13,11}{⼼,⾗}
  \begin{Phonetics}{想象}{xiang3xiang4}[][HSK 4]
    \definition[个,种,面]{s.}{imaginação; refere"-se ao processo mental de processamento e transformação de representações armazenadas na mente para formar novas imagens}
    \definition{v.}{imaginar; vislumbrar; visualizar; refere"-se a ter uma imagem concreta de algo que não está na frente dos olhos}
  \end{Phonetics}
\end{Entry}

\begin{Entry}{想想看}{13,13,9}{⼼,⼼,⽬}
  \begin{Phonetics}{想想看}{xiang3xiang3kan4}
    \definition{v.}{pensar sobre isso}
  \end{Phonetics}
\end{Entry}

%%%%%%%%%% 愁 %%%%%%%%%%
\subsection*{愁}\addcontentsline{loh}{figure}{愁}

\begin{Entry}{愁}{13}{⼼}
  \begin{Phonetics}{愁}{chou2}[][HSK 5]
    \definition{adj.}{triste; pesaroso; angustiado; desconsolado; preocupado; deprimido}
    \definition{s.}{pesar; sofrimento; dor; tristeza}
    \definition{v.}{preocupar"-se; estar preocupado; ficar ansioso; sentir ansiedade}
  \end{Phonetics}
\end{Entry}

\begin{Entry}{愁眉苦脸}{13,9,8,11}{⼼,⽬,⾋,⾁}
  \begin{Phonetics}{愁眉苦脸}{chou2mei2-ku3lian3}[][HSK 7-9]
    \definition{expr.}{com uma expressão preocupada e angustiada; parecendo perturbado; usar uma expressão preocupada; fazer uma cara feia}
  \end{Phonetics}
\end{Entry}

%%%%%%%%%% 愈 %%%%%%%%%%
\subsection*{愈}\addcontentsline{loh}{figure}{愈}

\begin{Entry}{愈}{13}{⼼}
  \begin{Phonetics}{愈}{yu4}
    \definition{v.}{curar; recuperar; ficar bem | Literário: ser melhor que; destacar"-se}
    \definition{v.}{(em 愈\dots 愈\dots) quanto mais\dots quanto mais\dots | mais; cada vez mais; a repetição indica que o grau se desenvolve à medida que as condições se desenvolvem}
  \seealsoref{愈……愈……}{yu4 yu4}
  \end{Phonetics}
\end{Entry}

\begin{Entry}{愈合}{13,6}{⼼,⼝}
  \begin{Phonetics}{愈合}{yu4he2}[][HSK 7-9]
    \definition{v.}{curar (feridas ou lesões cicatrizadas)}
  \synonymref{治愈}{zhi4yu4}
  \antonymref{创伤}{chuang1shang1}
  \antonymref{分裂}{fen1lie4}
  \antonymref{伤痕}{shang1hen2}
  \end{Phonetics}
\end{Entry}

\begin{Entry}{愈来愈}{13,7,13}{⼼,⽊,⼼}
  \begin{Phonetics}{愈来愈}{yu4lai2yu4}[][HSK 7-9]
    \definition{expr.}{cada vez mais; mais e mais}
  \end{Phonetics}
\end{Entry}

\begin{Entry}{愈……愈……}{13,13}{⼼,⼼}
  \begin{Phonetics}{愈……愈……}{yu4 yu4}
    \definition{conj.}{quanto mais\dots quanto mais\dots}
  \seealsoref{愈}{yu4}
  \end{Phonetics}
\end{Entry}

\begin{Entry}{愈演愈烈}{13,14,13,10}{⼼,⽔,⼼,⽕}
  \begin{Phonetics}{愈演愈烈}{yu4yan3-yu4lie4}[][HSK 7-9]
    \definition{expr.}{escalando; aumentar em intensidade; tornar"-se cada vez mais intenso; piorar; cada vez mais crítico; os problemas se tornam cada vez mais intensos}
  \end{Phonetics}
\end{Entry}

%%%%%%%%%% 意 %%%%%%%%%%
\subsection*{意}\addcontentsline{loh}{figure}{意}

\begin{Entry}{意}{13}{⼼}
  \begin{Phonetics}{意}{yi4}
    \definition*{s.}{Itália, abreviação de 意大利}
    \definition{s.}{ideia; significado; pensamento | desejo; vontade; intenção | significância}
  \seealsoref{意大利}{yi4da4li4}
  \end{Phonetics}
\end{Entry}

\begin{Entry}{意义}{13,3}{⼼,⼂}
  \begin{Phonetics}{意义}{yi4yi4}[][HSK 3]
    \definition[个,种,层,重,点]{s.}{sentido; significado; o significado expresso por meio de linguagem escrita ou outros sinais; o significado identificado por meio de ações ou obtenção | valor; efeito; significado; influência; impacto}
  \end{Phonetics}
\end{Entry}

\begin{Entry}{意大利}{13,3,7}{⼼,⼤,⼑}
  \begin{Phonetics}{意大利}{yi4da4li4}
    \definition*{s.}{Itália}
  \end{Phonetics}
\end{Entry}

\begin{Entry}{意见}{13,4}{⼼,⾒}
  \begin{Phonetics}{意见}{yi4jian5}[][HSK 2]
    \definition[种,点,条]{s.}{ideia; visão; opinião; sugestão; uma certa visão ou ideia sobre algo | objeção; reclamação; opinião divergente; (em relação a uma pessoa ou coisa) o sentimento de estar insatisfeito com algo porque está errado}
  \synonymref{观点}{guan1dian3}
  \synonymref{见解}{jian4jie3}
  \synonymref{建议}{jian4yi4}
  \synonymref{看法}{kan4fa5}
  \synonymref{想法}{xiang3fa5}
  \synonymref{主张}{zhu3zhang1}
  \synonymref{主意}{zhu3yi5}
  \end{Phonetics}
\end{Entry}

\begin{Entry}{意外}{13,5}{⼼,⼣}
  \begin{Phonetics}{意外}{yi4wai4}[][HSK 3]
    \definition{adj.}{inesperado; imprevisto}
    \definition{adv.}{acidentalmente}
    \definition[个,种]{s.}{acidente; infortúnio; um infortúnio inesperado}
  \end{Phonetics}
\end{Entry}

\begin{Entry}{意向}{13,6}{⼼,⼝}
  \begin{Phonetics}{意向}{yi4xiang4}[][HSK 7-9]
    \definition{s.}{intenção; propósito; intenções e planos}
  \seealsoref{意图}{yi4tu2}
  \synonymref{抱负}{bao4fu4}
  \synonymref{动向}{dong4xiang4}
  \synonymref{理想}{li3xiang3}
  \synonymref{希望}{xi1wang4}
  \synonymref{愿望}{yuan4wang4}
  \synonymref{志气}{zhi4qi4}
  \synonymref{志愿}{zhi4yuan4}
  \end{Phonetics}
\end{Entry}

\begin{Entry}{意志}{13,7}{⼼,⼼}
  \begin{Phonetics}{意志}{yi4zhi4}[][HSK 5]
    \definition[个,股]{s.}{vontade; determinação; desejo; força de vontade; o estado psicológico produzido pela determinação de atingir um determinado objetivo, muitas vezes expresso por meio de linguagem e ações}
  \end{Phonetics}
\end{Entry}

\begin{Entry}{意识}{13,7}{⼼,⾔}
  \begin{Phonetics}{意识}{yi4shi5}[][HSK 5]
    \definition{s.}{consciência; percepção; grau de reconhecimento e importância atribuído a uma determinada questão}
    \definition{s.}{consciência; percepção; o reflexo da mente humana no mundo material objetivo é a soma de vários processos psicológicos, como sensação e pensamento | consciência; conscientização; o grau de conscientização e atenção dada a um problema}
    \definition{v.}{perceber; despertar para; estar ciente de; sentir, descobrir o que antes não se sentia ou não se descobria; geralmente é usado junto com 到}
  \seealsoref{到}{dao4}
  \end{Phonetics}
\end{Entry}

\begin{Entry}{意译}{13,7}{⼼,⾔}
  \begin{Phonetics}{意译}{yi4yi4}
    \definition{s.}{tradução liberal | tradução semântica | tradução livre | significado (de expressão estrangeira) | paráfrase | tradução do significado}
  \synonymref{翻译}{fan1yi4}
  \antonymref{直译}{zhi2yi4}
  \end{Phonetics}
\end{Entry}

\begin{Entry}{意味着}{13,8,11}{⼼,⼝,⽬}
  \begin{Phonetics}{意味着}{yi4wei4zhe5}[][HSK 5]
    \definition{v.}{significar; subentender; implicar; entender que tem vários significados}
  \end{Phonetics}
\end{Entry}

\begin{Entry}{意图}{13,8}{⼼,⼞}
  \begin{Phonetics}{意图}{yi4tu2}[][HSK 7-9]
    \definition[个,种]{s.}{intenção; propósito; objetivo; a intenção de alcançar um determinado objetivo}
    \definition{v.}{tentar; pretender}
  \seealsoref{意向}{yi4xiang4}
  \synonymref{企图}{qi3tu2}
  \synonymref{妄图}{wang4tu2}
  \synonymref{妄想}{wang4xiang3}
  \synonymref{意愿}{yi4yuan4}
  \antonymref{无意}{wu2yi4}
  \end{Phonetics}
\end{Entry}

\begin{Entry}{意思}{13,9}{⼼,⼼}
  \begin{Phonetics}{意思}{yi4si5}[][HSK 2]
    \definition[个]{s.}{ideia; significado; o significado da linguagem e das palavras; conteúdo ideológico | desejo; vontade; opiniões | um símbolo de afeto, apreciação, gratidão, etc. | dica; traço; sugestão; refere"-se principalmente ao afeto entre homens e mulheres | diversão; interesse}
    \definition{v.}{dar uma dica; demonstrar sua gratidão com presentes ou outros meios}
  \seealsoref{口气}{kou3qi4}
  \seealsoref{意义}{yi4yi4}
  \synonymref{涵义}{han2yi4}
  \synonymref{含义}{han2yi4}
  \synonymref{乐趣}{le4qu4}
  \synonymref{趣味}{qu4wei4}
  \synonymref{思想}{si1xiang3}
  \synonymref{意图}{yi4tu2}
  \synonymref{意义}{yi4yi4}
  \end{Phonetics}
\end{Entry}

\begin{Entry}{意指}{13,9}{⼼,⼿}
  \begin{Phonetics}{意指}{yi4zhi3}
    \definition{v.}{implicar; inferir; pretender | significar}
  \end{Phonetics}
\end{Entry}

\begin{Entry}{意料}{13,10}{⼼,⽃}
  \begin{Phonetics}{意料}{yi4liao4}[][HSK 7-9]
    \definition{v.}{esperar; antecipar; estimar; estimativa prévia de circunstâncias, resultados, etc.}
  \synonymref{预见}{yu4jian4}
  \synonymref{预料}{yu4liao4}
  \antonymref{不料}{bu2liao4}
  \antonymref{意外}{yi4wai4}
  \end{Phonetics}
\end{Entry}

\begin{Entry}{意料之外}{13,10,3,5}{⼼,⽃,⼂,⼣}
  \begin{Phonetics}{意料之外}{yi4liao4zhi1wai4}[][HSK 7-9]
    \definition{expr.}{contrário ao esperado; inesperado}
  \synonymref{出人意料}{chu1ren2yi4liao4}
  \synonymref{意想不到}{yi4xiang3bu2dao4}
  \end{Phonetics}
\end{Entry}

\begin{Entry}{意想不到}{13,13,4,8}{⼼,⼼,⼀,⼑}
  \begin{Phonetics}{意想不到}{yi4xiang3bu2dao4}[][HSK 6]
    \definition{expr.}{anteriormente inimaginável | inesperado}
  \end{Phonetics}
\end{Entry}

\begin{Entry}{意愿}{13,14}{⼼,⽕}
  \begin{Phonetics}{意愿}{yi4yuan4}[][HSK 6]
    \definition{s.}{desejo; aspiração; vontade}
  \end{Phonetics}
\end{Entry}

%%%%%%%%%% 愚 %%%%%%%%%%
\subsection*{愚}\addcontentsline{loh}{figure}{愚}

\begin{Entry}{愚}{13}{⼼}
  \begin{Phonetics}{愚}{yu2}
    \definition{adj.}{tolo; estúpido}
    \definition{pron.}{eu; termos humildes usados para se referir a si mesmo}
    \definition{v.}{fazer de bobo; enganar}
  \synonymref{笨}{ben4}
  \synonymref{蠢}{chun3}
  \synonymref{傻}{sha3}
  \antonymref{智}{zhi4}
  \end{Phonetics}
\end{Entry}

\begin{Entry}{愚公移山}{13,4,11,3}{⼼,⼋,⽲,⼭}
  \begin{Phonetics}{愚公移山}{yu2gong1-yi2shan1}[][HSK 7-9]
    \definition{expr.}{``O velho insensato move a montanha.''; figurativamente, querer é poder; faça coisas aparentemente impossíveis com perseverança obstinada e, eventualmente, tenha sucesso; conta a lenda que, na antiguidade, havia um velho chamado Yu Gong, de Beishan; duas grandes montanhas bloqueavam seu caminho em frente à sua casa, e ele resolveu arrasá"-las; outro velho, Zhi Sou, de Hequ, riu de sua tolice, achando impossível. Yu Gong respondeu: ``Quando eu morrer, terei filhos; quando meus filhos morrerem, terei netos; nossa descendência é infinita. Essas duas montanhas não crescerão mais. Cada pedaço que removermos será um pouco a menos; um dia elas serão arrasadas.'' (De Liezi, ``Tang Wen''《列子·汤问》); esta história é uma analogia para a perseverança e a coragem diante das dificuldades}
  \seealsoref{列子}{lie4zi3}
  \end{Phonetics}
\end{Entry}

\begin{Entry}{愚昧}{13,9}{⼼,⽇}
  \begin{Phonetics}{愚昧}{yu2mei4}
    \definition{adj.}{ignorante; desinformado; falta de conhecimento; tolo e ignorante}
  \seealsoref{愚蠢}{yu2chun3}
  \synonymref{笨拙}{ben4zhuo1}
  \synonymref{痴呆}{chi1dai1}
  \synonymref{封建}{feng1jian4}
  \synonymref{迷信}{mi2xin4}
  \synonymref{无知}{wu2zhi1}
  \synonymref{愚蠢}{yu2chun3}
  \antonymref{聪明}{cong1ming5}
  \antonymref{机智}{ji1zhi4}
  \antonymref{机灵}{ji1ling5}
  \antonymref{英明}{ying1ming2}
  \antonymref{智慧}{zhi4hui4}
  \end{Phonetics}
\end{Entry}

\begin{Entry}{愚蠢}{13,21}{⼼,⾍}
  \begin{Phonetics}{愚蠢}{yu2chun3}[][HSK 7-9]
    \definition{adj.}{estúpido; tolo; bobo; não inteligente}
  \seealsoref{笨拙}{ben4zhuo1}
  \seealsoref{愚昧}{yu2mei4}
  \synonymref{笨拙}{ben4zhuo1}
  \synonymref{痴呆}{chi1dai1}
  \synonymref{无知}{wu2zhi1}
  \synonymref{愚昧}{yu2mei4}
  \antonymref{聪慧}{cong1hui4}
  \antonymref{聪明}{cong1ming5}
  \antonymref{明智}{ming2zhi4}
  \antonymref{巧妙}{qiao3miao4}
  \antonymref{智慧}{zhi4hui4}
  \end{Phonetics}
\end{Entry}

%%%%%%%%%% 感 %%%%%%%%%%
\subsection*{感}\addcontentsline{loh}{figure}{感}

\begin{Entry}{感}{13}{⼼}
  \begin{Phonetics}{感}{gan3}[][HSK 7-9]
    \definition{s.}{sentido; sensação; sentimento; impressão | emoção; sentimento}
    \definition{v.}{sentir; perceber; estar ciente | mover; tocar; afetar | ser grato; ser agradecido | ser afetado (pelo frio); pegar um resfriado | (fotografia) sensibilizar | ser grato; apreciar | ser afetado}
  \end{Phonetics}
\end{Entry}

\begin{Entry}{感人}{13,2}{⼼,⼈}
  \begin{Phonetics}{感人}{gan3ren2}[][HSK 6]
    \definition{adj.}{comovente; tocante}
  \end{Phonetics}
\end{Entry}

\begin{Entry}{感叹}{13,5}{⼼,⼝}
  \begin{Phonetics}{感叹}{gan3tan4}[][HSK 7-9]
    \definition{v.}{suspirar com sentimento; suspirar por causa de um sentimento}
  \end{Phonetics}
\end{Entry}

\begin{Entry}{感兴趣}{13,6,15}{⼼,⼋,⾛}
  \begin{Phonetics}{感兴趣}{gan3/xing4qu4}[][HSK 4]
    \definition{v.+compl.}{estar interessado}
  \seealsoref{对……感兴趣}{dui4 gan3xing4qu4}
  \end{Phonetics}
\end{Entry}

\begin{Entry}{感动}{13,6}{⼼,⼒}
  \begin{Phonetics}{感动}{gan3dong4}[][HSK 2]
    \definition{v.}{mover (alguém) | tocar (alguém emocionalmente)}
  \seealsoref{激动}{ji1dong4}
  \synonymref{触动}{chu4dong4}
  \synonymref{打动}{da3dong4}
  \synonymref{感触}{gan3chu4}
  \synonymref{感激}{gan3ji1}
  \synonymref{感人}{gan3ren2}
  \synonymref{激动}{ji1dong4}
  \antonymref{冷淡}{leng3dan4}
  \antonymref{冷漠}{leng3mo4}
  \antonymref{漠然}{mo4ran2}
  \end{Phonetics}
\end{Entry}

\begin{Entry}{感到}{13,8}{⼼,⼑}
  \begin{Phonetics}{感到}{gan3dao4}[][HSK 2]
    \definition{v.}{sentir; achar; perceber}
  \seealsoref{感觉}{gan3jue2}
  \seealsoref{觉得}{jue2de5}
  \synonymref{感触}{gan3chu4}
  \synonymref{感觉}{gan3jue2}
  \synonymref{觉得}{jue2de5}
  \synonymref{受到}{shou4dao4}
  \antonymref{无视}{wu2shi4}
  \end{Phonetics}
\end{Entry}

\begin{Entry}{感受}{13,8}{⼼,⼜}
  \begin{Phonetics}{感受}{gan3shou4}[][HSK 3]
    \definition{s.}{percepção ; compreenção; sentimento; experiência; influência do contato com o mundo exterior}
    \definition{v.}{sentir; sentir (através dos sentidos); experimentar; ser afetado}
  \seealsoref{感觉}{gan3jue2}
  \synonymref{感触}{gan3chu4}
  \synonymref{感觉}{gan3jue2}
  \synonymref{感想}{gan3xiang3}
  \synonymref{体会}{ti3hui4}
  \antonymref{抗拒}{kang4ju4}
  \antonymref{无视}{wu2shi4}
  \end{Phonetics}
\end{Entry}

\begin{Entry}{感性}{13,8}{⼼,⼼}
  \begin{Phonetics}{感性}{gan3xing4}[][HSK 7-9]
    \definition{adj.}{perceptivo; sentimental; emocional; pertencente a formas intuitivas como sensação, percepção e representação}
    \definition{s.}{percepção; sensibilidade; refere"-se a pessoas que são emocionalmente ricas, sentimentais, capazes de ter empatia pelos outros, que têm grande sensibilidade e conseguem entender as mudanças emocionais de qualquer coisa}
  \end{Phonetics}
\end{Entry}

\begin{Entry}{感冒}{13,9}{⼼,⽇}
  \begin{Phonetics}{感冒}{gan3mao4}[][HSK 3]
    \definition{adj.}{interessado em}
    \definition[场,次]{s.}{resfriado; gripe comum; \emph{influenza}; doença infecciosa causada por um vírus, que tende a causar sintomas como garganta seca, congestão nasal, tosse, espirros, dor de cabeça e febre quando o corpo está excessivamente cansado, resfriado ou com a imunidade enfraquecida}
    \definition{v.}{pegar (ter) um resfriado}
  \synonymref{着凉}{zhao2liang2}
  \end{Phonetics}
\end{Entry}

\begin{Entry}{感染}{13,9}{⼼,⽊}
  \begin{Phonetics}{感染}{gan3ran3}[][HSK 7-9]
    \definition{v.}{infectar; ser infectado com | infectar; afetar; influenciar; evocar os mesmos pensamentos e sentimentos}
  \end{Phonetics}
\end{Entry}

\begin{Entry}{感染力}{13,9,2}{⼼,⽊,⼒}
  \begin{Phonetics}{感染力}{gan3ran3li4}[][HSK 7-9]
    \definition{s.}{poder de mover os sentimentos; apelo | infeccioso (entusiasmo) | inspiração}
  \end{Phonetics}
\end{Entry}

\begin{Entry}{感觉}{13,9}{⼼,⾒}
  \begin{Phonetics}{感觉}{gan3jue2}[][HSK 2]
    \definition[个]{s.}{sentimento; sensação; percepção sensorial}
    \definition{v.}{sentir; perceber; tomar consciência; sentir no coração, acreditar}
  \seealsoref{感到}{gan3dao4}
  \seealsoref{感受}{gan3shou4}
  \seealsoref{觉得}{jue2de5}
  \synonymref{发觉}{fa1jue2}
  \synonymref{感触}{gan3chu4}
  \synonymref{感到}{gan3dao4}
  \synonymref{感受}{gan3shou4}
  \synonymref{感想}{gan3xiang3}
  \synonymref{觉得}{jue2de5}
  \synonymref{体会}{ti3hui4}
  \synonymref{嗅觉}{xiu4jue2}
  \antonymref{明确}{ming2que4}
  \antonymref{确定}{que4ding4}
  \end{Phonetics}
\end{Entry}

\begin{Entry}{感恩}{13,10}{⼼,⼼}
  \begin{Phonetics}{感恩}{gan3/en1}[][HSK 7-9]
    \definition{v.+compl.}{sentir"-se grato; ser grato; expressar gratidão pela ajuda dada por outros}
  \end{Phonetics}
\end{Entry}

\begin{Entry}{感情}{13,11}{⼼,⼼}
  \begin{Phonetics}{感情}{gan3qing2}[][HSK 3]
    \definition[份,个,种]{s.}{emoção; sentimento; reações psicológicas como amor, ódio, alegria, raiva, tristeza e felicidade, geradas por estímulos externos | amor; afeto; apego; preocupação e afeição por pessoas ou coisas}
  \seealsoref{情感}{qing2gan3}
  \synonymref{情感}{qing2gan3}
  \synonymref{情绪}{qing2xu4}
  \synonymref{心情}{xin1qing2}
  \end{Phonetics}
\end{Entry}

\begin{Entry}{感慨}{13,12}{⼼,⼼}
  \begin{Phonetics}{感慨}{gan3kai3}[][HSK 7-9]
    \definition{v.}{suspirar de emoção; estar cheio de emoções; ficar profundamente comovido;  geralmente expresso em palavras}
  \end{Phonetics}
\end{Entry}

\begin{Entry}{感谢}{13,12}{⼼,⾔}
  \begin{Phonetics}{感谢}{gan3xie4}[][HSK 2]
    \definition{v.}{agradecer; ser grato; expressar gratidão com palavras ou ações}
  \seealsoref{感激}{gan3ji1}
  \seealsoref{谢谢}{xie4xie5}
  \synonymref{报答}{bao4da2}
  \synonymref{感激}{gan3ji1}
  \synonymref{谢谢}{xie4xie5}
  \antonymref{抱怨}{bao4yuan5}
  \end{Phonetics}
\end{Entry}

\begin{Entry}{感想}{13,13}{⼼,⼼}
  \begin{Phonetics}{感想}{gan3xiang3}[][HSK 5]
    \definition[个,条]{s.}{pensamentos; impressões; reflexões; resposta do pensamento decorrente da exposição ao mundo exterior}
  \end{Phonetics}
\end{Entry}

\begin{Entry}{感触}{13,13}{⼼,⾓}
  \begin{Phonetics}{感触}{gan3chu4}[][HSK 7-9]
    \definition{s.}{sentimento; pensamentos e sentimentos; pensamentos e emoções causados por fatores externos}
  \end{Phonetics}
\end{Entry}

\begin{Entry}{感激}{13,16}{⼼,⽔}
  \begin{Phonetics}{感激}{gan3ji1}[][HSK 7-9]
    \definition{v.}{apreciar; ser grato; sentir"-se grato; sentir"-se em dívida; desenvolver uma impressão favorável de alguém por causa de sua gentileza ou ajuda}
  \end{Phonetics}
\end{Entry}

%%%%%%%%%% 慈 %%%%%%%%%%
\subsection*{慈}\addcontentsline{loh}{figure}{慈}

\begin{Entry}{慈}{13}{⼼}
  \begin{Phonetics}{慈}{ci2}
    \definition*{s.}{Sobrenome: Ci}
    \definition{adj.}{gentil; amoroso}
    \definition{s.}{mãe; refere"-se à mãe}
    \definition{v.}{Literário: amar (amor de cima para baixo)}
  \end{Phonetics}
\end{Entry}

\begin{Entry}{慈祥}{13,10}{⼼,⽰}
  \begin{Phonetics}{慈祥}{ci2xiang2}[][HSK 7-9]
    \definition{adj.}{gentil; benigno; amável; descreve a aparência ou atitude de uma pessoa idosa como sendo gentil, amável e acessível}
  \end{Phonetics}
\end{Entry}

\begin{Entry}{慈善}{13,12}{⼼,⼝}
  \begin{Phonetics}{慈善}{ci2shan4}[][HSK 7-9]
    \definition{adj.}{caridoso; benevolente; filantrópico; gentil e caridoso}
  \end{Phonetics}
\end{Entry}

%%%%%%%%%% 慎 %%%%%%%%%%
\subsection*{慎}\addcontentsline{loh}{figure}{慎}

\begin{Entry}{慎}{13}{⼼}
  \begin{Phonetics}{慎}{shen4}
    \definition*{s.}{Sobrenome: Shen}
    \definition{adj.}{cuidadoso; cauteloso}
  \end{Phonetics}
\end{Entry}

\begin{Entry}{慎重}{13,9}{⼼,⾥}
  \begin{Phonetics}{慎重}{shen4zhong4}[][HSK 7-9]
    \definition{adj.}{cuidadoso; cauteloso; discreto; prudente}
  \synonymref{把稳}{ba3 wen3}
  \synonymref{谨慎}{jin3shen4}
  \synonymref{留心}{liu2/xin1}
  \synonymref{留意}{liu2/yi4}
  \synonymref{小心}{xiao3xin5}
  \end{Phonetics}
\end{Entry}

%%%%%%%%%% 慢 %%%%%%%%%%
\subsection*{慢}\addcontentsline{loh}{figure}{慢}

\begin{Entry}{慢}{14}{⼼}
  \begin{Phonetics}{慢}{man4}[][HSK 1]
    \definition*{s.}{Sobrenome: Man}
    \definition{adj.}{lento; devagar; baixa velocidade; longa duração | rude; arrogante; sem educação com as pessoas | frouxo; lento}
    \definition{adv.}{lentamente}
  \synonymref{迟}{chi2}
  \synonymref{缓}{huan3}
  \antonymref{急}{ji2}
  \antonymref{快}{kuai4}
  \end{Phonetics}
\end{Entry}

\begin{Entry}{慢车}{14,4}{⼼,⾞}
  \begin{Phonetics}{慢车}{man4che1}[][HSK 6]
    \definition{s.}{trem lento com muitas paradas | ônibus ou trem local; parada do trem}
  \antonymref{快车}{kuai4che1}
  \end{Phonetics}
\end{Entry}

\begin{Entry}{慢动作}{14,6,7}{⼼,⼒,⼈}
  \begin{Phonetics}{慢动作}{man4dong4zuo4}
    \definition{s.}{micromovimento; ação lenta; movimento lento | câmera lenta}
  \end{Phonetics}
\end{Entry}

\begin{Entry}{慢性}{14,8}{⼼,⼼}
  \begin{Phonetics}{慢性}{man4xing4}[][HSK 7-9]
    \definition{adj.}{crônico; duradouro | lento (em fazer efeito)}
  \end{Phonetics}
\end{Entry}

\begin{Entry}{慢慢}{14,14}{⼼,⼼}
  \begin{Phonetics}{慢慢}{man4man4}[][HSK 3]
    \definition{adv.}{lentamente; vagarosamente; gradualmente | lentamente; vagarosamente; gradualmente; depois de um longo período de tempo}
  \end{Phonetics}
\end{Entry}

\begin{Entry}{慢慢来}{14,14,7}{⼼,⼼,⽊}
  \begin{Phonetics}{慢慢来}{man4man4lai2}[][HSK 7-9]
    \definition{v.}{ir com calma; não ter pressa; significa não ter impaciência ao fazer as coisas e prosseguir no seu próprio ritmo}
  \end{Phonetics}
\end{Entry}

%%%%%%%%%% 慷 %%%%%%%%%%
\subsection*{慷}\addcontentsline{loh}{figure}{慷}

\begin{Entry}{慷}{14}{⼼}
  \begin{Phonetics}{慷}{kang1}
    \definition{adj.}{generoso | magnânimo}
  \end{Phonetics}
\end{Entry}

\begin{Entry}{慷慨}{14,12}{⼼,⼼}
  \begin{Phonetics}{慷慨}{kang1kai3}[][HSK 7-9]
    \definition{adj.}{generoso; descreve alguém como alguém que não é mesquinho; disposto a ajudar os outros com dinheiro ou bens | veemente; fervoroso; apaixonado; descreve alguém como alguém repleto de senso de justiça e emocionalmente intenso}
  \end{Phonetics}
\end{Entry}

%%%%%%%%%% 慰 %%%%%%%%%%
\subsection*{慰}\addcontentsline{loh}{figure}{慰}

\begin{Entry}{慰}{15}{⼼}
  \begin{Phonetics}{慰}{wei4}
    \definition{adj.}{aliviado; em paz; confortável}
    \definition{v.}{consolar; confortar | ser (ficar) aliviado}
  \end{Phonetics}
\end{Entry}

\begin{Entry}{慰问}{15,6}{⼼,⾨}
  \begin{Phonetics}{慰问}{wei4wen4}[][HSK 5]
    \definition{v.}{visitar; consolar; expressar simpatia por; confortar e cumprimentar com palavras e presentes;  enfatizar o conforto e o cumprimento, frequentemente usado por superiores para subordinados}
  \end{Phonetics}
\end{Entry}

\begin{Entry}{慰劳}{15,7}{⼼,⼒}
  \begin{Phonetics}{慰劳}{wei4lao2}[][HSK 7-9]
    \definition{v.}{levar presentes ou enviar votos de felicidades em reconhecimento aos serviços prestados | confortar | demonstrar apreço (através de palavras gentis, pequenos presentes, etc.)}
  \synonymref{安慰}{an1wei4}
  \synonymref{慰问}{wei4wen4}
  \synonymref{问候}{wen4hou4}
  \antonymref{打击}{da3ji1}
  \end{Phonetics}
\end{Entry}

%%%%%%%%%% 憋 %%%%%%%%%%
\subsection*{憋}\addcontentsline{loh}{figure}{憋}

\begin{Entry}{憋}{15}{⼼}
  \begin{Phonetics}{憋}{bie1}[][HSK 7-9]
    \definition{adj.}{sufocado; oprimido}
    \definition{v.}{suprimir; conter | Dialeto: obrigar | Dialeto: ponderar; contemplar | Dialeto: ficar de olho em | Dialeto: destruir (por pressão interna) | calar a boca; inibir; bloquear | sufocar; abafar}
  \end{Phonetics}
\end{Entry}

%%%%%%%%%% 憧 %%%%%%%%%%
\subsection*{憧}\addcontentsline{loh}{figure}{憧}

\begin{Entry}{憧}{15}{⼼}
  \begin{Phonetics}{憧}{chong1}
    \definition{adj.}{irresoluto; indeciso | estúpido; imbecil; confuso}
  \end{Phonetics}
\end{Entry}

\begin{Entry}{憧憬}{15,15}{⼼,⼼}
  \begin{Phonetics}{憧憬}{chong1jing3}
    \definition{v.}{ansiar por; desejar ardentemente; esperar com expectativa}
  \synonymref{期待}{qi1dai4}
  \synonymref{期望}{qi1wang4}
  \synonymref{希望}{xi1wang4}
  \synonymref{向往}{xiang4wang3}
  \synonymref{展望}{zhan3wang4}
  \antonymref{绝望}{jue2/wang4}
  \antonymref{失望}{shi1wang4}
  \antonymref{遗忘}{yi2wang4}
  \end{Phonetics}
\end{Entry}

%%%%%%%%%% 懂 %%%%%%%%%%
\subsection*{懂}\addcontentsline{loh}{figure}{懂}

\begin{Entry}{懂}{15}{⼼}
  \begin{Phonetics}{懂}{dong3}[][HSK 2]
    \definition*{s.}{Sobrenome: Dong}
    \definition{v.}{compreender; entender}
  \seealsoref{明白}{ming2bai5}
  \seealsoref{知道}{zhi1dao5}
  \synonymref{明白}{ming2bai5}
  \synonymref{晓}{xiao3}
  \synonymref{知}{zhi1}
  \synonymref{知道}{zhi1dao5}
  \antonymref{懵}{meng3}
  \antonymref{无知}{wu2zhi1}
  \end{Phonetics}
\end{Entry}

\begin{Entry}{懂事}{15,8}{⼼,⼅}
  \begin{Phonetics}{懂事}{dong3shi4}[][HSK 7-9]
    \definition{adj.}{sensato; inteligente; muito compreensivo da natureza e da razão humana}
  \end{Phonetics}
\end{Entry}

\begin{Entry}{懂得}{15,11}{⼼,⼻}
  \begin{Phonetics}{懂得}{dong3de5}[][HSK 2]
    \definition{v.}{saber (significado, prática, etc.); compreender; entender}
  \seealsoref{晓得}{xiao3de2}
  \seealsoref{知道}{zhi1dao5}
  \synonymref{理会}{li3hui4}
  \synonymref{理解}{li3jie3}
  \synonymref{了解}{liao3jie3}
  \synonymref{领会}{ling3hui4}
  \synonymref{清楚}{qing1chu5}
  \synonymref{学会}{xue2hui4}
  \synonymref{知道}{zhi1dao5}
  \end{Phonetics}
\end{Entry}

%%%%%%%%%% 懒 %%%%%%%%%%
\subsection*{懒}\addcontentsline{loh}{figure}{懒}

\begin{Entry}{懒}{16}{⼼}
  \begin{Phonetics}{懒}{lan3}[][HSK 6]
    \definition{adj.}{indolente; preguiçoso | lento; lânguido | ocioso; preguiçoso}
  \antonymref{勤}{qin2}
  \end{Phonetics}
\end{Entry}

\begin{Entry}{懒人}{16,2}{⼼,⼈}
  \begin{Phonetics}{懒人}{lan3ren2}
    \definition{s.}{pessoa preguiçosa}
  \synonymref{懒虫}{lan3chong2}
  \synonymref{懒汉}{lan3han4}
  \synonymref{懒散}{lan3san3}
  \end{Phonetics}
\end{Entry}

\begin{Entry}{懒汉}{16,5}{⼼,⽔}
  \begin{Phonetics}{懒汉}{lan3han4}
    \definition{s.}{pessoa preguiçosa; pessoa ociosa; pessoa indolente}
  \synonymref{懒人}{lan3ren2}
  \end{Phonetics}
\end{Entry}

\begin{Entry}{懒虫}{16,6}{⼼,⾍}
  \begin{Phonetics}{懒虫}{lan3chong2}
    \definition{s.}{pessoa preguiçosa (um insulto ou um termo humorístico)}
  \synonymref{懒鬼}{lan3gui3}
  \synonymref{懒汉}{lan3han4}
  \synonymref{懒人}{lan3ren2}
  \end{Phonetics}
\end{Entry}

\begin{Entry}{懒怠}{16,9}{⼼,⼼}
  \begin{Phonetics}{懒怠}{lan3dai4}
    \definition{adj.}{preguiçoso; indolente}
    \definition{adv.}{preguiçoso demais para; indisposto a | não estou com vontade de}
  \synonymref{不屑}{bu2xie4}
  \synonymref{懒惰}{lan3duo4}
  \synonymref{懒散}{lan3san3}
  \synonymref{冷淡}{leng3dan4}
  \synonymref{无意}{wu2yi4}
  \antonymref{积极}{ji1ji2}
  \antonymref{乐意}{le4yi4}
  \antonymref{勤奋}{qin2fen4}
  \antonymref{勤劳}{qin2lao2}
  \antonymref{勤快}{qin2kuai5}
  \antonymref{热衷}{re4zhong1}
  \antonymref{愿意}{yuan4yi5}
  \end{Phonetics}
\end{Entry}

\begin{Entry}{懒鬼}{16,9}{⼼,⿁}
  \begin{Phonetics}{懒鬼}{lan3gui3}
    \definition{s.}{pessoa preguiçosa; pessoa indolente; vadio; vagabundo}
  \synonymref{懒虫}{lan3chong2}
  \antonymref{勤快}{qin2kuai5}
  \end{Phonetics}
\end{Entry}

\begin{Entry}{懒得}{16,11}{⼼,⼻}
  \begin{Phonetics}{懒得}{lan3de5}[][HSK 7-9]
    \definition{v.}{não estar disposto a fazer algo; estar entediado; estar sem vontade (de fazer algo)}
  \end{Phonetics}
\end{Entry}

\begin{Entry}{懒惰}{16,12}{⼼,⼼}
  \begin{Phonetics}{懒惰}{lan3duo4}[][HSK 7-9]
    \definition{adj.}{preguiçoso; ocioso; que não gosta de trabalho e esforço; que não é diligente}
  \end{Phonetics}
\end{Entry}

\begin{Entry}{懒散}{16,12}{⼼,⽁}
  \begin{Phonetics}{懒散}{lan3san3}
    \definition{adj.}{apático; relaxado; lento; descreve uma pessoa como apática, indisciplinada e sem motivação}
  \synonymref{懒怠}{lan3dai4}
  \synonymref{懒惰}{lan3duo4}
  \synonymref{偷懒}{tou1/lan3}
  \antonymref{勤奋}{qin2fen4}
  \antonymref{勤劳}{qin2lao2}
  \antonymref{勤快}{qin2kuai5}
  \end{Phonetics}
\end{Entry}

\begin{Entry}{懒腰}{16,13}{⼼,⾁}
  \begin{Phonetics}{懒腰}{lan3yao1}
    \definition[个]{v.}{alongar"-se; refere"-se a um alongamento preguiçoso da cintura ou das costas, geralmente indicando fadiga ou o ato de acordar}
  \end{Phonetics}
\end{Entry}

%%%%%%%%%% 懵 %%%%%%%%%%
\subsection*{懵}\addcontentsline{loh}{figure}{懵}

\begin{Entry}{懵}{18}{⼼}
  \begin{Phonetics}{懵}{meng3}
    \definition{adj.}{confuso; ignorante; irracional | inconsciente; entorpecido}
  \end{Phonetics}
\end{Entry}

\begin{Entry}{懵懂}{18,15}{⼼,⼼}
  \begin{Phonetics}{懵懂}{meng3dong3}
    \definition{adj.}{confuso; ignorante dos fatos}
  \synonymref{糊涂}{hu2tu5}
  \antonymref{明白}{ming2bai5}
  \antonymref{醒目}{xing3mu4}
  \end{Phonetics}
\end{Entry}

%%%%%%%%%% 聼 %%%%%%%%%%
\subsection*{聼}\addcontentsline{loh}{figure}{聼}

\begin{Entry}{聼}{19}{⼼}
  \begin{Phonetics}{聼}{ting1}
    \variantof{听}
  \end{Phonetics}
\end{Entry}

%%%%% EOF %%%%%

